% Complete reviewed Unit 7; editorial master: chapters/ch07/student.md
% Markdown SHA-256: ed4887966801964ff9eeab2b90ee7da65ba49edb598a57e125e35089cf4d3046
% One-use initial assembly followed by complete content and page review.
\newcommand{\runningtitle}{Unit 7. Subjunctive and passive}
% Layout inherited from the reviewed previous course; set a running title per unit.
% This is a small hand-maintained TeX layout, not a conversion pipeline.
\RequirePackage{fix-cm}
\documentclass[11pt,twoside]{article}
\usepackage{fontspec}
\setmainfont{Cambria}[
  Renderer=HarfBuzz,
  BoldFont=cambriab.ttf,
  ItalicFont=cambriai.ttf,
  BoldItalicFont=cambriaz.ttf
]
\setmonofont{lmmono10-regular.otf}[BoldFont=lmmonolt10-bold.otf]
\usepackage[paperwidth=165mm,paperheight=241mm,inner=21mm,outer=19mm,
  top=22mm,bottom=20mm,headheight=14pt,headsep=7mm,footskip=10mm]{geometry}
\usepackage{xcolor,array,booktabs,longtable,tabularx}
\usepackage{enumitem,fancyhdr}
\usepackage[unicode,hidelinks]{hyperref}
\definecolor{ink}{HTML}{000000}
\definecolor{muted}{HTML}{4C5558}
\setlength{\parindent}{0pt}
\setlength{\parskip}{3.5pt}
\AtBeginDocument{\fontsize{11}{14.3}\selectfont}
\setlength{\tabcolsep}{4pt}
\renewcommand{\arraystretch}{1.10}
\setlist[itemize]{leftmargin=1.3em,label=\textbullet,itemsep=2pt,topsep=3pt,parsep=0pt}
\setlist[enumerate]{leftmargin=1.8em,itemsep=3pt,topsep=4pt,parsep=0pt}
\widowpenalty=10000
\clubpenalty=10000
\emergencystretch=1em
\raggedbottom
\pagestyle{fancy}
\fancyhf{}
\providecommand{\runningtitle}{Old English}
\fancyhead[LE,RO]{\small\itshape \runningtitle}
\fancyhead[LO,RE]{\small Old English}
\fancyfoot[LE,RO]{\small\thepage}
\renewcommand{\headrulewidth}{0.3pt}
\newcommand{\booktitle}[2]{
  {\fontsize{9}{11}\selectfont\color{muted}#1\par}
  \vspace{3pt}
  {\fontsize{20}{24}\selectfont\color{ink}#2\par}
  \vspace{6pt}
}
\newcommand{\source}[1]{\par{\footnotesize\color{muted}#1\par}}
\newcommand{\readingtext}[1]{
  \par\addvspace{5pt}
  \begingroup\leftskip=1.5em
    \fontsize{12}{16}\selectfont #1\par
  \endgroup\addvspace{5pt}
}
\makeatletter
\renewcommand{\section}{\@startsection{section}{1}{\z@}
  {-2.4ex \@plus -0.6ex \@minus -0.2ex}{0.8ex \@plus 0.2ex}
  {\normalfont\fontsize{12}{15}\selectfont\bfseries\color{ink}}}
\renewcommand{\subsection}{\@startsection{subsection}{2}{\z@}
  {-1.8ex \@plus -0.5ex \@minus -0.2ex}{0.5ex \@plus 0.2ex}
  {\normalfont\fontsize{11}{14}\selectfont\bfseries\color{ink}}}
\makeatother

\hyphenpenalty=10000
\exhyphenpenalty=10000
\emergencystretch=3em
\newcommand{\itempara}[2]{\par\addvspace{3pt}\noindent\hangindent=1.8em\hangafter=1\makebox[1.8em][l]{#1.}#2\par}
\newcommand{\reserve}[1]{\par\penalty-100\begingroup\dimen0=\pagegoal\advance\dimen0 by-\pagetotal\ifdim\dimen0<#1\newpage\fi\endgroup}
\begin{document}
\addfontfeatures{Ligatures=TeXOff}
\raggedright
\booktitle{UNIT 7}{Subjunctive, participles, and passive}

% mdparagraph:P001:begin
This unit adds wishes and exhortations, basic conditional clauses, participles, and passive constructions. It gives complete present and preterite subjunctives for the verbs named below, together with the new vocabulary needed for Cook I.6–7.
% mdparagraph:P001:end

% mdparagraph:P002:begin
In verbal tables, \textbf{1}, \textbf{2}, and \textbf{3} mean first, second, and third person; \textbf{sg.} and \textbf{pl.} mean singular and plural. A row headed \textbf{1–3 sg.} or \textbf{1–3 pl.} includes all three persons. Dual subjects take plural finite forms. Noun and adjective tables use nominative, genitive, dative, and accusative; the shorter labels are \textbf{nom}, \textbf{gen}, \textbf{dat}, and \textbf{acc}. Preserve vowel length and the distinction between e and ę. Examples and the account “Binding the animals” are new teaching material; the final Creation paragraph is an original source excerpt.
% mdparagraph:P002:end

\reserve{7\baselineskip}
% mdparagraph:P003:begin
\section*{1. Assertions, wishes, and commands}
% mdparagraph:P003:end

% teaching: NEW-U07-GRAM-U07-WISH-EXHORTATION
% mdparagraph:P004:begin
The \textbf{indicative} presents an assertion. The \textbf{subjunctive}, called the \textbf{optative} by Cook, can express a wish or an exhortation. In an independent clause its present tense can mean “may …” or “let …”. This English “let” translates the force of the inflected verb; no separate Old English verb meaning “let” is needed.
% mdparagraph:P004:end

% example: NEW-U07-EXAMPLE-01
% mdparagraph:P005:begin
Hīe swimmað on ðǣm sǣ. Swimmen hīe on ðǣm sǣ.
% mdparagraph:P005:end

% mdparagraph:P006:begin
They swim in the sea. Let them swim in the sea.
% mdparagraph:P006:end

% mdparagraph:P007:begin
The first clause asserts an activity; the second calls for it. Swimmað is indicative; swimmen is subjunctive. Both have the plural subject hīe. The sea phrase uses the masculine sǣ paradigm learned earlier.
% mdparagraph:P007:end

% example: NEW-U07-EXAMPLE-02
% mdparagraph:P008:begin
Gāð. Gān wē.
% mdparagraph:P008:end

% mdparagraph:P009:begin
Go. Let us go.
% mdparagraph:P009:end

% mdparagraph:P010:begin
Gāð is an imperative addressed to several people. Gān with wē is a first-person plural subjunctive exhortation: it includes the speaker. The selected subjunctive plural retains -n before its following pronoun.
% mdparagraph:P010:end

% mdparagraph:P011:begin
The familiar imperative gives a direct second-person command. The subjunctive permits first- and third-person wishes as well. The person belongs to the subject and context: its ending distinguishes number, but not first, second, and third person within that number.
% mdparagraph:P011:end

% teaching: NEW-U07-GRAM-U07-MODAL-MOOD-DISTINCTION
% example: NEW-U07-EXAMPLE-03
\reserve{4\baselineskip}
% mdparagraph:P012:begin
Hē mæg swimman. Hē mōt swimman. Swimme hē.
% mdparagraph:P012:end

% mdparagraph:P013:begin
He can swim. He may swim. Let him swim.
% mdparagraph:P013:end

% mdparagraph:P014:begin
Mæg asserts capability; mōt asserts permission. Each governs the plain infinitive swimman. Swimme expresses a wish about his swimming without asserting either capability or permission.
% mdparagraph:P014:end

% mdparagraph:P015:begin
Keep the earlier meanings of cunnan “know how”, ðurfan “need”, durran “dare”, sculan “be obliged”, and willan “wish, want” distinct from inflectional mood. A statement that someone wants something can itself be indicative. Willan is not automatically an English future auxiliary, and wolde is not automatically “would”. Where indicative and subjunctive spellings coincide, the clause setting must decide the intended use.
% mdparagraph:P015:end

\reserve{7\baselineskip}
% mdparagraph:P016:begin
\section*{2. The complete subjunctive systems}
% mdparagraph:P016:end

% teaching: NEW-U07-GRAM-U07-SUBJUNCTIVE-ENDINGS
% teaching: NEW-U07-GRAM-U07-SUBJUNCTIVE-STEMS
% mdparagraph:P017:begin
The ordinary endings are \textbf{-e} throughout the singular and \textbf{-en} throughout the plural, in both present and preterite. Thus four displayed forms for each verb give all twelve person–number–tense combinations. Contracted and irregular forms have their actual spellings in the tables.
% mdparagraph:P017:end

% mdparagraph:P018:begin
For the present, use the selected present stem. Do not carry the indicative second- or third-singular vowel mutation into the whole subjunctive. For a strong preterite, use the \textbf{preterite plural stem}, with its vowel and consonant changes, in both numbers. A weak preterite retains its dental stem. The preterite indicative plural -on and subjunctive plural -en remain distinct.
% mdparagraph:P018:end

% mdparagraph:P019:begin
The named verbs define this unit’s complete subjunctive assignment. They do not assign a subjunctive system to every verb previously encountered in a reading.
% mdparagraph:P019:end

\reserve{7\baselineskip}
% mdparagraph:P020:begin
\subsection*{Weak class I}
% mdparagraph:P020:end

% teaching: NEW-U07-SUBJ-hieran
% teaching: NEW-U07-SUBJ-sellan
% teaching: NEW-U07-SUBJ-gewyrcean
% teaching: NEW-U07-SUBJ-wyrcean
% teaching: NEW-U07-SUBJ-gefremman
% teaching: NEW-U07-SUBJ-gierwan
% teaching: NEW-U07-SUBJ-nerian
% mdparagraph:P021:begin
Keep the irregular past stems with their verbs: sęllan has sealde, and gewyrcean/\allowbreak{}wyrcean have geworhte/\allowbreak{}worhte. Gefręmman retains double mm in the present, but its past has single m before -ede. Gierwan retains w in the selected present and has gierede in the past. Nęrian retains i in -ie/\allowbreak{}-ien; its past is -ede, despite the infinitive’s resemblance to weak class II.
% mdparagraph:P021:end

% table: U07-DISPLAY-SUBJ-1
% mdtable:T01:begin
% course-table-style: 2026-10
% course-panel:T01-P01:begin
\par\begingroup
\fontsize{10.5}{13.2}\selectfont
\setlength{\LTleft}{0pt}\setlength{\LTright}{\fill}
\noindent
\begin{tabular}{@{}>{\raggedright\arraybackslash}p{20.873mm}>{\raggedright\arraybackslash}p{21.064mm}>{\raggedright\arraybackslash}p{24.876mm}>{\raggedright\arraybackslash}p{21.064mm}>{\raggedright\arraybackslash}p{24.876mm}@{}}
\toprule
\textbf{Person and number} & \textbf{hīeran · Present} & \textbf{hīeran · Preterite} & \textbf{sęllan · Present} & \textbf{sęllan · Preterite} \\
\midrule
1–3 sg. & hīere & hīerde & sęlle & sealde \\
1–3 pl. & hīeren & hīerden & sęllen & sealden \\
\bottomrule
\end{tabular}
\par\endgroup
% course-panel:T01-P01:end
\par\addvspace{5pt}
% course-panel:T01-P02:begin
\par\begingroup
\fontsize{10.5}{13.2}\selectfont
\setlength{\LTleft}{0pt}\setlength{\LTright}{\fill}
\noindent
\begin{tabular}{@{}>{\raggedright\arraybackslash}p{19.336mm}>{\raggedright\arraybackslash}p{23.839mm}>{\raggedright\arraybackslash}p{25.429mm}>{\raggedright\arraybackslash}p{20.308mm}>{\raggedright\arraybackslash}p{23.839mm}@{}}
\toprule
\textbf{Person and number} & \textbf{gewyrcean · Present} & \textbf{gewyrcean · Preterite} & \textbf{wyrcean · Present} & \textbf{wyrcean · Preterite} \\
\midrule
1–3 sg. & gewyrce & geworhte & wyrce & worhte \\
1–3 pl. & gewyrcen & geworhten & wyrcen & worhten \\
\bottomrule
\end{tabular}
\par\endgroup
% course-panel:T01-P02:end
\par\addvspace{5pt}
% course-panel:T01-P03:begin
\par\begingroup
\fontsize{10.5}{13.2}\selectfont
\setlength{\LTleft}{0pt}\setlength{\LTright}{\fill}
\noindent
\begin{tabular}{@{}>{\raggedright\arraybackslash}p{19.171mm}>{\raggedright\arraybackslash}p{23.636mm}>{\raggedright\arraybackslash}p{26.175mm}>{\raggedright\arraybackslash}p{20.134mm}>{\raggedright\arraybackslash}p{23.636mm}@{}}
\toprule
\textbf{Person and number} & \textbf{gefręmman · Present} & \textbf{gefręmman · Preterite} & \textbf{gierwan · Present} & \textbf{gierwan · Preterite} \\
\midrule
1–3 sg. & gefręmme & gefręmede & gierwe & gierede \\
1–3 pl. & gefręmmen & gefręmeden & gierwen & giereden \\
\bottomrule
\end{tabular}
\par\endgroup
% course-panel:T01-P03:end
\par\addvspace{5pt}
\par\addvspace{3pt}
% course-caption:T01:3:begin
\textbf{nęrian}\par\nobreak
% course-caption:T01:3:end
% course-panel:T01-P04:begin
\par\begingroup
\fontsize{10.5}{13.2}\selectfont
\setlength{\LTleft}{0pt}\setlength{\LTright}{\fill}
\noindent
\begin{tabular}{@{}>{\raggedright\arraybackslash}p{39.459mm}>{\raggedright\arraybackslash}p{39.459mm}>{\raggedright\arraybackslash}p{39.459mm}@{}}
\toprule
\textbf{Person and number} & \textbf{Present} & \textbf{Preterite} \\
\midrule
1–3 sg. & nęrie & nęrede \\
1–3 pl. & nęrien & nęreden \\
\bottomrule
\end{tabular}
\par\endgroup
% course-panel:T01-P04:end
% mdtable:T01:end
\reserve{7\baselineskip}
% mdparagraph:P022:begin
\subsection*{Weak classes II and III}
% mdparagraph:P022:end

% teaching: NEW-U07-SUBJ-lufian
% teaching: NEW-U07-SUBJ-gebletsian
% teaching: NEW-U07-SUBJ-gehalgian
% teaching: NEW-U07-SUBJ-habban
% teaching: NEW-U07-SUBJ-libban
% teaching: NEW-U07-SUBJ-secgan
% teaching: NEW-U07-SUBJ-hycgan
% mdparagraph:P023:begin
The first three verbs are class II: present -ie/\allowbreak{}-ien, preterite -ode/\allowbreak{}-oden. The last four are class III and require their own stems. Habban has hæbb- in the present subjunctive; its indicative plural does not supply that vowel. Libban has lifd- in the past; sęcgan has sægd-; hycgan has hogod-.
% mdparagraph:P023:end

% table: U07-DISPLAY-SUBJ-2
% mdtable:T02:begin
% course-table-style: 2026-10
% course-panel:T02-P01:begin
\par\begingroup
\fontsize{10.5}{13.2}\selectfont
\setlength{\LTleft}{0pt}\setlength{\LTright}{\fill}
\noindent
\begin{tabular}{@{}>{\raggedright\arraybackslash}p{18.849mm}>{\raggedright\arraybackslash}p{19.022mm}>{\raggedright\arraybackslash}p{22.464mm}>{\raggedright\arraybackslash}p{24.96mm}>{\raggedright\arraybackslash}p{27.456mm}@{}}
\toprule
\textbf{Person and number} & \textbf{lufian · Present} & \textbf{lufian · Preterite} & \textbf{geblētsian · Present} & \textbf{geblētsian · Preterite} \\
\midrule
1–3 sg. & lufie & lufode & geblētsie & geblētsode \\
1–3 pl. & lufien & lufoden & geblētsien & geblētsoden \\
\bottomrule
\end{tabular}
\par\endgroup
% course-panel:T02-P01:end
\par\addvspace{5pt}
% course-panel:T02-P02:begin
\par\begingroup
\fontsize{10.5}{13.2}\selectfont
\setlength{\LTleft}{0pt}\setlength{\LTright}{\fill}
\noindent
\begin{tabular}{@{}>{\raggedright\arraybackslash}p{19.443mm}>{\raggedright\arraybackslash}p{23.971mm}>{\raggedright\arraybackslash}p{26.546mm}>{\raggedright\arraybackslash}p{19.621mm}>{\raggedright\arraybackslash}p{23.172mm}@{}}
\toprule
\textbf{Person and number} & \textbf{gehālgian · Present} & \textbf{gehālgian · Preterite} & \textbf{habban · Present} & \textbf{habban · Preterite} \\
\midrule
1–3 sg. & gehālgie & gehālgode & hæbbe & hæfde \\
1–3 pl. & gehālgien & gehālgoden & hæbben & hæfden \\
\bottomrule
\end{tabular}
\par\endgroup
% course-panel:T02-P02:end
\par\addvspace{5pt}
% course-panel:T02-P03:begin
\par\begingroup
\fontsize{10.5}{13.2}\selectfont
\setlength{\LTleft}{0pt}\setlength{\LTright}{\fill}
\noindent
\begin{tabular}{@{}>{\raggedright\arraybackslash}p{20.873mm}>{\raggedright\arraybackslash}p{21.064mm}>{\raggedright\arraybackslash}p{24.876mm}>{\raggedright\arraybackslash}p{21.064mm}>{\raggedright\arraybackslash}p{24.876mm}@{}}
\toprule
\textbf{Person and number} & \textbf{libban · Present} & \textbf{libban · Preterite} & \textbf{sęcgan · Present} & \textbf{sęcgan · Preterite} \\
\midrule
1–3 sg. & libbe & lifde & sęcge & sægde \\
1–3 pl. & libben & lifden & sęcgen & sægden \\
\bottomrule
\end{tabular}
\par\endgroup
% course-panel:T02-P03:end
\par\addvspace{5pt}
\par\addvspace{3pt}
% course-caption:T02:3:begin
\textbf{hycgan}\par\nobreak
% course-caption:T02:3:end
% course-panel:T02-P04:begin
\par\begingroup
\fontsize{10.5}{13.2}\selectfont
\setlength{\LTleft}{0pt}\setlength{\LTright}{\fill}
\noindent
\begin{tabular}{@{}>{\raggedright\arraybackslash}p{39.459mm}>{\raggedright\arraybackslash}p{39.459mm}>{\raggedright\arraybackslash}p{39.459mm}@{}}
\toprule
\textbf{Person and number} & \textbf{Present} & \textbf{Preterite} \\
\midrule
1–3 sg. & hycge & hogode \\
1–3 pl. & hycgen & hogoden \\
\bottomrule
\end{tabular}
\par\endgroup
% course-panel:T02-P04:end
% mdtable:T02:end
% mdparagraph:P024:begin
\textbf{Alternatives.} Libban also has present lifie/\allowbreak{}lifien; sęcgan has past sāde/\allowbreak{}sāden. The shorter past of hycgan beside the selected hogode-series is also documented. These alternatives do not replace the selected table forms in a task that expressly asks for them.
% mdparagraph:P024:end

\reserve{7\baselineskip}
% mdparagraph:P025:begin
\subsection*{Strong classes I–III}
% mdparagraph:P025:end

% teaching: NEW-U07-SUBJ-drifan
% teaching: NEW-U07-SUBJ-wreon
% teaching: NEW-U07-SUBJ-beodan
% teaching: NEW-U07-SUBJ-teon
% teaching: NEW-U07-SUBJ-fleogan
% teaching: NEW-U07-SUBJ-bindan
% teaching: NEW-U07-SUBJ-helpan
% teaching: NEW-U07-SUBJ-weorthan
% teaching: NEW-U07-SUBJ-swimman
% mdparagraph:P026:begin
Use the familiar plural-preterite stems, not the first-singular past stems. Bindan therefore has bunde/\allowbreak{}bunden, and helpan has hulpe/\allowbreak{}hulpen. In drīfan the present ī and past i remain distinct. Weorðan changes ð to d in the selected past. The contracted wrēon and tēon have singular wrēo and tēo, without an added -e; their past stems contain g.
% mdparagraph:P026:end

% table: U07-DISPLAY-SUBJ-3
% mdtable:T03:begin
% course-table-style: 2026-10
% course-panel:T03-P01:begin
\par\begingroup
\fontsize{10.5}{13.2}\selectfont
\setlength{\LTleft}{0pt}\setlength{\LTright}{\fill}
\noindent
\begin{tabular}{@{}>{\raggedright\arraybackslash}p{21.195mm}>{\raggedright\arraybackslash}p{21.389mm}>{\raggedright\arraybackslash}p{25.26mm}>{\raggedright\arraybackslash}p{20.518mm}>{\raggedright\arraybackslash}p{24.389mm}@{}}
\toprule
\textbf{Person and number} & \textbf{drīfan · Present} & \textbf{drīfan · Preterite} & \textbf{wrēon · Present} & \textbf{wrēon · Preterite} \\
\midrule
1–3 sg. & drīfe & drife & wrēo & wrige \\
1–3 pl. & drīfen & drifen & wrēon & wrigen \\
\bottomrule
\end{tabular}
\par\endgroup
% course-panel:T03-P01:end
\par\addvspace{5pt}
% course-panel:T03-P02:begin
\par\begingroup
\fontsize{10.5}{13.2}\selectfont
\setlength{\LTleft}{0pt}\setlength{\LTright}{\fill}
\noindent
\begin{tabular}{@{}>{\raggedright\arraybackslash}p{21.528mm}>{\raggedright\arraybackslash}p{21.725mm}>{\raggedright\arraybackslash}p{25.657mm}>{\raggedright\arraybackslash}p{19.955mm}>{\raggedright\arraybackslash}p{23.887mm}@{}}
\toprule
\textbf{Person and number} & \textbf{bēodan · Present} & \textbf{bēodan · Preterite} & \textbf{tēon · Present} & \textbf{tēon · Preterite} \\
\midrule
1–3 sg. & bēode & bude & tēo & tuge \\
1–3 pl. & bēoden & buden & tēon & tugen \\
\bottomrule
\end{tabular}
\par\endgroup
% course-panel:T03-P02:end
\par\addvspace{5pt}
% course-panel:T03-P03:begin
\par\begingroup
\fontsize{10.5}{13.2}\selectfont
\setlength{\LTleft}{0pt}\setlength{\LTright}{\fill}
\noindent
\begin{tabular}{@{}>{\raggedright\arraybackslash}p{20.56mm}>{\raggedright\arraybackslash}p{21.593mm}>{\raggedright\arraybackslash}p{25.348mm}>{\raggedright\arraybackslash}p{20.748mm}>{\raggedright\arraybackslash}p{24.503mm}@{}}
\toprule
\textbf{Person and number} & \textbf{flēogan · Present} & \textbf{flēogan · Preterite} & \textbf{bindan · Present} & \textbf{bindan · Preterite} \\
\midrule
1–3 sg. & flēoge & fluge & binde & bunde \\
1–3 pl. & flēogen & flugen & binden & bunden \\
\bottomrule
\end{tabular}
\par\endgroup
% course-panel:T03-P03:end
\par\addvspace{5pt}
% course-panel:T03-P04:begin
\par\begingroup
\fontsize{10.5}{13.2}\selectfont
\setlength{\LTleft}{0pt}\setlength{\LTright}{\fill}
\noindent
\begin{tabular}{@{}>{\raggedright\arraybackslash}p{20.56mm}>{\raggedright\arraybackslash}p{20.748mm}>{\raggedright\arraybackslash}p{24.503mm}>{\raggedright\arraybackslash}p{21.593mm}>{\raggedright\arraybackslash}p{25.348mm}@{}}
\toprule
\textbf{Person and number} & \textbf{helpan · Present} & \textbf{helpan · Preterite} & \textbf{weorðan · Present} & \textbf{weorðan · Preterite} \\
\midrule
1–3 sg. & helpe & hulpe & weorðe & wurde \\
1–3 pl. & helpen & hulpen & weorðen & wurden \\
\bottomrule
\end{tabular}
\par\endgroup
% course-panel:T03-P04:end
\par\addvspace{5pt}
\par\addvspace{3pt}
% course-caption:T03:4:begin
\textbf{swimman}\par\nobreak
% course-caption:T03:4:end
% course-panel:T03-P05:begin
\par\begingroup
\fontsize{10.5}{13.2}\selectfont
\setlength{\LTleft}{0pt}\setlength{\LTright}{\fill}
\noindent
\begin{tabular}{@{}>{\raggedright\arraybackslash}p{39.459mm}>{\raggedright\arraybackslash}p{39.459mm}>{\raggedright\arraybackslash}p{39.459mm}@{}}
\toprule
\textbf{Person and number} & \textbf{Present} & \textbf{Preterite} \\
\midrule
1–3 sg. & swimme & swumme \\
1–3 pl. & swimmen & swummen \\
\bottomrule
\end{tabular}
\par\endgroup
% course-panel:T03-P05:end
% mdtable:T03:end
% mdparagraph:P027:begin
Tēon here is the class II verb “draw, pull”. The same spelling can name another verb; it is the learned lexical system that determines the past stem.
% mdparagraph:P027:end

\reserve{7\baselineskip}
% mdparagraph:P028:begin
\subsection*{Strong classes IV–VII}
% mdparagraph:P028:end

% teaching: NEW-U07-SUBJ-beran
% teaching: NEW-U07-SUBJ-niman
% teaching: NEW-U07-SUBJ-cuman
% teaching: NEW-U07-SUBJ-giefan
% teaching: NEW-U07-SUBJ-cwethan
% teaching: NEW-U07-SUBJ-biddan
% teaching: NEW-U07-SUBJ-licgan
% teaching: NEW-U07-SUBJ-sittan
% teaching: NEW-U07-SUBJ-seon
% teaching: NEW-U07-SUBJ-faran
% teaching: NEW-U07-SUBJ-standan
% teaching: NEW-U07-SUBJ-slean
% teaching: NEW-U07-SUBJ-hebban
% teaching: NEW-U07-SUBJ-scieppan
% teaching: NEW-U07-SUBJ-swerian
% teaching: NEW-U07-SUBJ-weaxan
% teaching: NEW-U07-SUBJ-hatan
% teaching: NEW-U07-SUBJ-fon
% mdparagraph:P029:begin
Retain the selected past-plural vowels of niman and cuman. Cweðan has d in cwǣde/\allowbreak{}cwǣden. The j-present verbs retain their characteristic present consonants: dd, cg, tt, bb, and pp. Swęrian retains i in the present ending. Standan loses its present n in stōde/\allowbreak{}stōden. The contracted sēon, slēan, and fōn require their own present and past forms.
% mdparagraph:P029:end

% table: U07-DISPLAY-SUBJ-4
% mdtable:T04:begin
% course-table-style: 2026-10
% course-panel:T04-P01:begin
\par\begingroup
\fontsize{10.5}{13.2}\selectfont
\setlength{\LTleft}{0pt}\setlength{\LTright}{\fill}
\noindent
\begin{tabular}{@{}>{\raggedright\arraybackslash}p{21.528mm}>{\raggedright\arraybackslash}p{20.84mm}>{\raggedright\arraybackslash}p{24.772mm}>{\raggedright\arraybackslash}p{20.84mm}>{\raggedright\arraybackslash}p{24.772mm}@{}}
\toprule
\textbf{Person and number} & \textbf{beran · Present} & \textbf{beran · Preterite} & \textbf{niman · Present} & \textbf{niman · Preterite} \\
\midrule
1–3 sg. & bere & bǣre & nime & nōme \\
1–3 pl. & beren & bǣren & nimen & nōmen \\
\bottomrule
\end{tabular}
\par\endgroup
% course-panel:T04-P01:end
\par\addvspace{5pt}
% course-panel:T04-P02:begin
\par\begingroup
\fontsize{10.5}{13.2}\selectfont
\setlength{\LTleft}{0pt}\setlength{\LTright}{\fill}
\noindent
\begin{tabular}{@{}>{\raggedright\arraybackslash}p{21.195mm}>{\raggedright\arraybackslash}p{20.518mm}>{\raggedright\arraybackslash}p{24.389mm}>{\raggedright\arraybackslash}p{21.389mm}>{\raggedright\arraybackslash}p{25.26mm}@{}}
\toprule
\textbf{Person and number} & \textbf{cuman · Present} & \textbf{cuman · Preterite} & \textbf{giefan · Present} & \textbf{giefan · Preterite} \\
\midrule
1–3 sg. & cume & cōme & giefe & gēafe \\
1–3 pl. & cumen & cōmen & giefen & gēafen \\
\bottomrule
\end{tabular}
\par\endgroup
% course-panel:T04-P02:end
\par\addvspace{5pt}
% course-panel:T04-P03:begin
\par\begingroup
\fontsize{10.5}{13.2}\selectfont
\setlength{\LTleft}{0pt}\setlength{\LTright}{\fill}
\noindent
\begin{tabular}{@{}>{\raggedright\arraybackslash}p{20.873mm}>{\raggedright\arraybackslash}p{21.064mm}>{\raggedright\arraybackslash}p{24.876mm}>{\raggedright\arraybackslash}p{21.064mm}>{\raggedright\arraybackslash}p{24.876mm}@{}}
\toprule
\textbf{Person and number} & \textbf{cweðan · Present} & \textbf{cweðan · Preterite} & \textbf{biddan · Present} & \textbf{biddan · Preterite} \\
\midrule
1–3 sg. & cweðe & cwǣde & bidde & bǣde \\
1–3 pl. & cweðen & cwǣden & bidden & bǣden \\
\bottomrule
\end{tabular}
\par\endgroup
% course-panel:T04-P03:end
\par\addvspace{5pt}
% course-panel:T04-P04:begin
\par\begingroup
\fontsize{10.5}{13.2}\selectfont
\setlength{\LTleft}{0pt}\setlength{\LTright}{\fill}
\noindent
\begin{tabular}{@{}>{\raggedright\arraybackslash}p{20.873mm}>{\raggedright\arraybackslash}p{21.064mm}>{\raggedright\arraybackslash}p{24.876mm}>{\raggedright\arraybackslash}p{21.064mm}>{\raggedright\arraybackslash}p{24.876mm}@{}}
\toprule
\textbf{Person and number} & \textbf{licgan · Present} & \textbf{licgan · Preterite} & \textbf{sittan · Present} & \textbf{sittan · Preterite} \\
\midrule
1–3 sg. & licge & lǣge & sitte & sǣte \\
1–3 pl. & licgen & lǣgen & sitten & sǣten \\
\bottomrule
\end{tabular}
\par\endgroup
% course-panel:T04-P04:end
\par\addvspace{5pt}
% course-panel:T04-P05:begin
\par\begingroup
\fontsize{10.5}{13.2}\selectfont
\setlength{\LTleft}{0pt}\setlength{\LTright}{\fill}
\noindent
\begin{tabular}{@{}>{\raggedright\arraybackslash}p{21.871mm}>{\raggedright\arraybackslash}p{20.273mm}>{\raggedright\arraybackslash}p{24.268mm}>{\raggedright\arraybackslash}p{21.172mm}>{\raggedright\arraybackslash}p{25.167mm}@{}}
\toprule
\textbf{Person and number} & \textbf{sēon · Present} & \textbf{sēon · Preterite} & \textbf{faran · Present} & \textbf{faran · Preterite} \\
\midrule
1–3 sg. & sēo & sāwe & fare & fōre \\
1–3 pl. & sēon & sāwen & faren & fōren \\
\bottomrule
\end{tabular}
\par\endgroup
% course-panel:T04-P05:end
\par\addvspace{5pt}
% course-panel:T04-P06:begin
\par\begingroup
\fontsize{10.5}{13.2}\selectfont
\setlength{\LTleft}{0pt}\setlength{\LTright}{\fill}
\noindent
\begin{tabular}{@{}>{\raggedright\arraybackslash}p{20.873mm}>{\raggedright\arraybackslash}p{21.921mm}>{\raggedright\arraybackslash}p{25.734mm}>{\raggedright\arraybackslash}p{20.206mm}>{\raggedright\arraybackslash}p{24.018mm}@{}}
\toprule
\textbf{Person and number} & \textbf{standan · Present} & \textbf{standan · Preterite} & \textbf{slēan · Present} & \textbf{slēan · Preterite} \\
\midrule
1–3 sg. & stande & stōde & slēa & slōge \\
1–3 pl. & standen & stōden & slēan & slōgen \\
\bottomrule
\end{tabular}
\par\endgroup
% course-panel:T04-P06:end
\par\addvspace{5pt}
% course-panel:T04-P07:begin
\par\begingroup
\fontsize{10.5}{13.2}\selectfont
\setlength{\LTleft}{0pt}\setlength{\LTright}{\fill}
\noindent
\begin{tabular}{@{}>{\raggedright\arraybackslash}p{20.075mm}>{\raggedright\arraybackslash}p{20.259mm}>{\raggedright\arraybackslash}p{23.925mm}>{\raggedright\arraybackslash}p{22.917mm}>{\raggedright\arraybackslash}p{25.575mm}@{}}
\toprule
\textbf{Person and number} & \textbf{hębban · Present} & \textbf{hębban · Preterite} & \textbf{scieppan · Present} & \textbf{scieppan · Preterite} \\
\midrule
1–3 sg. & hębbe & hōfe & scieppe & scēope \\
1–3 pl. & hębben & hōfen & scieppen & scēopen \\
\bottomrule
\end{tabular}
\par\endgroup
% course-panel:T04-P07:end
\par\addvspace{5pt}
% course-panel:T04-P08:begin
\par\begingroup
\fontsize{10.5}{13.2}\selectfont
\setlength{\LTleft}{0pt}\setlength{\LTright}{\fill}
\noindent
\begin{tabular}{@{}>{\raggedright\arraybackslash}p{20.56mm}>{\raggedright\arraybackslash}p{21.593mm}>{\raggedright\arraybackslash}p{25.348mm}>{\raggedright\arraybackslash}p{20.748mm}>{\raggedright\arraybackslash}p{24.503mm}@{}}
\toprule
\textbf{Person and number} & \textbf{swęrian · Present} & \textbf{swęrian · Preterite} & \textbf{weaxan · Present} & \textbf{weaxan · Preterite} \\
\midrule
1–3 sg. & swęrie & swōre & weaxe & wēoxe \\
1–3 pl. & swęrien & swōren & weaxen & wēoxen \\
\bottomrule
\end{tabular}
\par\endgroup
% course-panel:T04-P08:end
\par\addvspace{5pt}
% course-panel:T04-P09:begin
\par\begingroup
\fontsize{10.5}{13.2}\selectfont
\setlength{\LTleft}{0pt}\setlength{\LTright}{\fill}
\noindent
\begin{tabular}{@{}>{\raggedright\arraybackslash}p{22.106mm}>{\raggedright\arraybackslash}p{21.4mm}>{\raggedright\arraybackslash}p{25.437mm}>{\raggedright\arraybackslash}p{20.188mm}>{\raggedright\arraybackslash}p{23.62mm}@{}}
\toprule
\textbf{Person and number} & \textbf{hātan · Present} & \textbf{hātan · Preterite} & \textbf{fōn · Present} & \textbf{fōn · Preterite} \\
\midrule
1–3 sg. & hāte & hēte & fō & fēnge \\
1–3 pl. & hāten & hēten & fōn & fēngen \\
\bottomrule
\end{tabular}
\par\endgroup
% course-panel:T04-P09:end
% mdtable:T04:end
% mdparagraph:P030:begin
Cuman also has the documented present cyme/\allowbreak{}cymen; the selected forms are cume/\allowbreak{}cumen. Keep scieppan’s selected scēop- past. Slēan retains the ordinary meaning “strike” learned earlier. Hāten is present plural subjunctive in its finite paradigm; the same spelling will occur as a separately identified past participle. Spelling alone cannot establish its function.
% mdparagraph:P030:end

\reserve{7\baselineskip}
% mdparagraph:P031:begin
\subsection*{Preterite-presents and anomalous verbs}
% mdparagraph:P031:end

% teaching: NEW-U07-SUBJ-witan
% teaching: NEW-U07-SUBJ-cunnan
% teaching: NEW-U07-SUBJ-thurfan
% teaching: NEW-U07-SUBJ-durran
% teaching: NEW-U07-SUBJ-sculan
% teaching: NEW-U07-SUBJ-magan
% teaching: NEW-U07-SUBJ-motan
% teaching: NEW-U07-SUBJ-willan
% teaching: NEW-U07-SUBJ-don
% teaching: NEW-U07-SUBJ-gan
% teaching: NEW-U07-SUBJ-wesan
% mdparagraph:P032:begin
The preterite-present verbs form their present subjunctives from the selected plural-type or mutated stems, not from a strong-looking indicative singular. Witan has wit-; ðurfan, durran, and sculan have the selected y forms; mugan has mæge. Their dental preterites take the ordinary singular/\allowbreak{}plural subjunctive contrast.
% mdparagraph:P032:end

% mdparagraph:P033:begin
Dōn and gān retain contracted dō/\allowbreak{}dōn and gā/\allowbreak{}gān in the present. Their preterites use dyde and ēode. Willan has wille/\allowbreak{}willen and wolde/\allowbreak{}wolden. Wesan has sīe/\allowbreak{}sīen and wǣre/\allowbreak{}wǣren.
% mdparagraph:P033:end

% table: U07-DISPLAY-SUBJ-5
% mdtable:T05:begin
% course-table-style: 2026-10
% course-panel:T05-P01:begin
\par\begingroup
\fontsize{10.5}{13.2}\selectfont
\setlength{\LTleft}{0pt}\setlength{\LTright}{\fill}
\noindent
\begin{tabular}{@{}>{\raggedright\arraybackslash}p{21.195mm}>{\raggedright\arraybackslash}p{20.518mm}>{\raggedright\arraybackslash}p{24.389mm}>{\raggedright\arraybackslash}p{21.389mm}>{\raggedright\arraybackslash}p{25.26mm}@{}}
\toprule
\textbf{Person and number} & \textbf{witan · Present} & \textbf{witan · Preterite} & \textbf{cunnan · Present} & \textbf{cunnan · Preterite} \\
\midrule
1–3 sg. & wite & wiste & cunne & cūðe \\
1–3 pl. & witen & wisten & cunnen & cūðen \\
\bottomrule
\end{tabular}
\par\endgroup
% course-panel:T05-P01:end
\par\addvspace{5pt}
% course-panel:T05-P02:begin
\par\begingroup
\fontsize{10.5}{13.2}\selectfont
\setlength{\LTleft}{0pt}\setlength{\LTright}{\fill}
\noindent
\begin{tabular}{@{}>{\raggedright\arraybackslash}p{20.873mm}>{\raggedright\arraybackslash}p{21.064mm}>{\raggedright\arraybackslash}p{24.876mm}>{\raggedright\arraybackslash}p{21.064mm}>{\raggedright\arraybackslash}p{24.876mm}@{}}
\toprule
\textbf{Person and number} & \textbf{ðurfan · Present} & \textbf{ðurfan · Preterite} & \textbf{durran · Present} & \textbf{durran · Preterite} \\
\midrule
1–3 sg. & ðyrfe & ðorfte & dyrre & dorste \\
1–3 pl. & ðyrfen & ðorften & dyrren & dorsten \\
\bottomrule
\end{tabular}
\par\endgroup
% course-panel:T05-P02:end
\par\addvspace{5pt}
% course-panel:T05-P03:begin
\par\begingroup
\fontsize{10.5}{13.2}\selectfont
\setlength{\LTleft}{0pt}\setlength{\LTright}{\fill}
\noindent
\begin{tabular}{@{}>{\raggedright\arraybackslash}p{21.195mm}>{\raggedright\arraybackslash}p{21.389mm}>{\raggedright\arraybackslash}p{25.26mm}>{\raggedright\arraybackslash}p{20.518mm}>{\raggedright\arraybackslash}p{24.389mm}@{}}
\toprule
\textbf{Person and number} & \textbf{sculan · Present} & \textbf{sculan · Preterite} & \textbf{mugan · Present} & \textbf{mugan · Preterite} \\
\midrule
1–3 sg. & scyle & sceolde & mæge & meahte \\
1–3 pl. & scylen & sceolden & mægen & meahten \\
\bottomrule
\end{tabular}
\par\endgroup
% course-panel:T05-P03:end
\par\addvspace{5pt}
% course-panel:T05-P04:begin
\par\begingroup
\fontsize{10.5}{13.2}\selectfont
\setlength{\LTleft}{0pt}\setlength{\LTright}{\fill}
\noindent
\begin{tabular}{@{}>{\raggedright\arraybackslash}p{21.195mm}>{\raggedright\arraybackslash}p{20.518mm}>{\raggedright\arraybackslash}p{24.389mm}>{\raggedright\arraybackslash}p{21.389mm}>{\raggedright\arraybackslash}p{25.26mm}@{}}
\toprule
\textbf{Person and number} & \textbf{mōtan · Present} & \textbf{mōtan · Preterite} & \textbf{willan · Present} & \textbf{willan · Preterite} \\
\midrule
1–3 sg. & mōte & mōste & wille & wolde \\
1–3 pl. & mōten & mōsten & willen & wolden \\
\bottomrule
\end{tabular}
\par\endgroup
% course-panel:T05-P04:end
\par\addvspace{5pt}
% course-panel:T05-P05:begin
\par\begingroup
\fontsize{10.5}{13.2}\selectfont
\setlength{\LTleft}{0pt}\setlength{\LTright}{\fill}
\noindent
\begin{tabular}{@{}>{\raggedright\arraybackslash}p{22.716mm}>{\raggedright\arraybackslash}p{20.746mm}>{\raggedright\arraybackslash}p{24.272mm}>{\raggedright\arraybackslash}p{20.746mm}>{\raggedright\arraybackslash}p{24.272mm}@{}}
\toprule
\textbf{Person and number} & \textbf{dōn · Present} & \textbf{dōn · Preterite} & \textbf{gān · Present} & \textbf{gān · Preterite} \\
\midrule
1–3 sg. & dō & dyde & gā & ēode \\
1–3 pl. & dōn & dyden & gān & ēoden \\
\bottomrule
\end{tabular}
\par\endgroup
% course-panel:T05-P05:end
\par\addvspace{5pt}
\par\addvspace{3pt}
% course-caption:T05:5:begin
\textbf{wesan}\par\nobreak
% course-caption:T05:5:end
% course-panel:T05-P06:begin
\par\begingroup
\fontsize{10.5}{13.2}\selectfont
\setlength{\LTleft}{0pt}\setlength{\LTright}{\fill}
\noindent
\begin{tabular}{@{}>{\raggedright\arraybackslash}p{39.459mm}>{\raggedright\arraybackslash}p{39.459mm}>{\raggedright\arraybackslash}p{39.459mm}@{}}
\toprule
\textbf{Person and number} & \textbf{Present} & \textbf{Preterite} \\
\midrule
1–3 sg. & sīe & wǣre \\
1–3 pl. & sīen & wǣren \\
\bottomrule
\end{tabular}
\par\endgroup
% course-panel:T05-P06:end
% mdtable:T05:end
% mdparagraph:P034:begin
\textbf{Alternative present of “be”.} Each row again covers all three persons. The bēo series shares the same preterite wǣre/\allowbreak{}wǣren; it has no separate past paradigm.
% mdparagraph:P034:end

% table: U07-DISPLAY-BEO-SUBJ
% mdtable:T06:begin
% course-table-style: 2026-10
% course-panel:T06-P01:begin
\par\begingroup
\fontsize{10.5}{13.2}\selectfont
\setlength{\LTleft}{0pt}\setlength{\LTright}{\fill}
\noindent
\begin{tabular}{@{}>{\raggedright\arraybackslash}p{63.342mm}>{\raggedright\arraybackslash}p{57.846mm}@{}}
\toprule
\textbf{Person and number} & \textbf{Present} \\
\midrule
1–3 sg. & bēo \\
1–3 pl. & bēon \\
\bottomrule
\end{tabular}
\par\endgroup
% course-panel:T06-P01:end
% mdtable:T06:end
% mdparagraph:P035:begin
\textbf{Alternatives and limits.} Witan has the alternative wisse/\allowbreak{}wissen preterite; cunnan has a cȳðe variant. Beside the selected ðyrfe/\allowbreak{}ðyrfen, dyrre/\allowbreak{}dyrren, and scyle/\allowbreak{}scylen, the corresponding forms with u are documented. The tables continue to select meahte/\allowbreak{}meahten for mugan. Learn the selected series together; a spelling variant is not a new mood.
% mdparagraph:P035:end

% mdparagraph:P036:begin
The labels mugan (also called magan), durran, sculan, and mōtan identify finite systems. Their presence in the table headings does not supply a usable infinitive, imperative, or participle. The new subjunctives do not fill every missing nonfinite or command form. Bēo itself can also be first-singular indicative or second-singular imperative; bēon can be an infinitive. A clause’s subject and purpose matter.
% mdparagraph:P036:end

\reserve{7\baselineskip}
% mdparagraph:P037:begin
\section*{3. Five additional verb systems}
% mdparagraph:P037:end

% teaching: NEW-U07-VOC-wyrcean-yield
% teaching: NEW-U07-VOC-gieman-1
% teaching: NEW-U07-VOC-ricsian-dominion
% teaching: NEW-U07-VOC-gemanigfieldan-multiply
% teaching: NEW-U07-VOC-totwaeman-separate
% teaching: NEW-U07-GRAM-U07-GIEMAN-GENITIVE-COMPLEMENT
% teaching: NEW-U07-GRAM-U07-RICSIAN-INTRANSITIVE
% teaching: NEW-U07-GRAM-U07-WYRCEAN-ACCUSATIVE-FRUIT
% teaching: NEW-U07-GRAM-U07-TOTWAEMAN-ACCUSATIVE-OBJECT-FRAM-DATIVE
% teaching: NEW-U07-GRAM-U07-GEMANIGFIELDAN-PASSIVE-NOMINATIVE-SUBJECT
% mdparagraph:P038:begin
These verbs use familiar weak patterns, but their meanings and complements must be learned individually. The five headwords identify the complete finite systems printed here. Only tōtwǣman adds a new imperative pair; the participle assignment is stated in §5.
% mdparagraph:P038:end

% mdtable:T07:begin
\par\begingroup
\fontsize{10.5}{13.2}\selectfont
\setlength{\LTleft}{0pt}\setlength{\LTright}{\fill}
\noindent
\begin{tabular}{@{}>{\raggedright\arraybackslash}p{26mm}>{\raggedright\arraybackslash}p{23mm}>{\raggedright\arraybackslash}p{69.37mm}@{}}
\toprule
\textbf{Verb} & \textbf{Class} & \textbf{Meaning and construction} \\
\midrule
wyrcean & weak I, irregular & yield fruit; the fruit is accusative \\
gīeman & weak I & rule over; the governed person or thing is genitive \\
rīcsian & weak II & have dominion, bear rule; used intransitively \\
gemanigfieldan & weak I, dental stem & multiply; this unit uses its passive “be multiplied” with the beings as subject \\
tōtwǣman & weak I & divide, separate; the thing divided is accusative, and what it is divided from follows fram in the dative \\
\bottomrule
\end{tabular}
\par\endgroup

% mdtable:T07:end
% mdparagraph:P039:begin
Wyrcean is taught here as its own verb. Knowing prefixed gewyrcean does not by itself establish an unprefixed meaning. The government of gīeman differs from an ordinary accusative object; rīcsian needs no object. Neither construction is an invitation to apply the accusative-object passive rule to it.
% mdparagraph:P039:end

% teaching: NEW-U07-GRAM-U07-FIVE-NEW-INDICATIVES
\reserve{7\baselineskip}
% mdparagraph:P040:begin
\subsection*{Present and preterite indicatives}
% mdparagraph:P040:end

% mdparagraph:P041:begin
The tables give each of the six person–number forms in both tenses. In the weak preterite, first and third singular coincide; second singular adds -st, and the plural has -on. Preserve the lexical stem and the dental changes shown.
% mdparagraph:P041:end

\reserve{5\baselineskip}
% mdparagraph:P042:begin
\textbf{Wyrcean.}
% mdparagraph:P042:end

% teaching: NEW-U07-IND-wyrcean
% mdparagraph:P043:begin
The present has wyrce-, with wyrcst and wyrcð in the second and third singular. The irregular past is worht-, not a regular addition to wyrc-.
% mdparagraph:P043:end

% table: NEW-U07-IND-wyrcean
% mdtable:T08:begin
% course-table-style: 2026-10
% course-panel:T08-P01:begin
\par\begingroup
\fontsize{10.5}{13.2}\selectfont
\setlength{\LTleft}{0pt}\setlength{\LTright}{\fill}
\noindent
\begin{tabular}{@{}>{\raggedright\arraybackslash}p{39.459mm}>{\raggedright\arraybackslash}p{39.459mm}>{\raggedright\arraybackslash}p{39.459mm}@{}}
\toprule
\textbf{Person and number} & \textbf{Present indicative} & \textbf{Preterite indicative} \\
\midrule
1 sg. & wyrce & worhte \\
2 sg. & wyrcst & worhtest \\
3 sg. & wyrcð & worhte \\
1–3 pl. & wyrceað & worhton \\
\bottomrule
\end{tabular}
\par\endgroup
% course-panel:T08-P01:end
% mdtable:T08:end
\reserve{5\baselineskip}
% mdparagraph:P044:begin
\textbf{Gīeman.}
% mdparagraph:P044:end

% teaching: NEW-U07-IND-gieman
% mdparagraph:P045:begin
The long stem has gīemst, gīemð, and plural gīemað; the preterite retains the dental stem gīemd-.
% mdparagraph:P045:end

% table: NEW-U07-IND-gieman
% mdtable:T09:begin
% course-table-style: 2026-10
% course-panel:T09-P01:begin
\par\begingroup
\fontsize{10.5}{13.2}\selectfont
\setlength{\LTleft}{0pt}\setlength{\LTright}{\fill}
\noindent
\begin{tabular}{@{}>{\raggedright\arraybackslash}p{39.459mm}>{\raggedright\arraybackslash}p{39.459mm}>{\raggedright\arraybackslash}p{39.459mm}@{}}
\toprule
\textbf{Person and number} & \textbf{Present indicative} & \textbf{Preterite indicative} \\
\midrule
1 sg. & gīeme & gīemde \\
2 sg. & gīemst & gīemdest \\
3 sg. & gīemð & gīemde \\
1–3 pl. & gīemað & gīemdon \\
\bottomrule
\end{tabular}
\par\endgroup
% course-panel:T09-P01:end
% mdtable:T09:end
\reserve{5\baselineskip}
% mdparagraph:P046:begin
\textbf{Rīcsian.}
% mdparagraph:P046:end

% teaching: NEW-U07-IND-ricsian
% mdparagraph:P047:begin
Class II contrasts rīcsie and plural rīcsiað with rīcsast/\allowbreak{}rīcsað. The preterite has -od-.
% mdparagraph:P047:end

% table: NEW-U07-IND-ricsian
% mdtable:T10:begin
% course-table-style: 2026-10
% course-panel:T10-P01:begin
\par\begingroup
\fontsize{10.5}{13.2}\selectfont
\setlength{\LTleft}{0pt}\setlength{\LTright}{\fill}
\noindent
\begin{tabular}{@{}>{\raggedright\arraybackslash}p{39.459mm}>{\raggedright\arraybackslash}p{39.459mm}>{\raggedright\arraybackslash}p{39.459mm}@{}}
\toprule
\textbf{Person and number} & \textbf{Present indicative} & \textbf{Preterite indicative} \\
\midrule
1 sg. & rīcsie & rīcsode \\
2 sg. & rīcsast & rīcsodest \\
3 sg. & rīcsað & rīcsode \\
1–3 pl. & rīcsiað & rīcsodon \\
\bottomrule
\end{tabular}
\par\endgroup
% course-panel:T10-P01:end
% mdtable:T10:end
\reserve{5\baselineskip}
% mdparagraph:P048:begin
\textbf{Gemanigfieldan.}
% mdparagraph:P048:end

% teaching: NEW-U07-IND-gemanigfieldan
% mdparagraph:P049:begin
The dental stem has the contracted present gemanigfieltst and gemanigfielt. The first-singular present and first-/\allowbreak{}third-singular preterite have the same spelling gemanigfielde; use the sentence to determine tense. Do not add another dental to the printed past stem.
% mdparagraph:P049:end

% table: NEW-U07-IND-gemanigfieldan
% mdtable:T11:begin
% course-table-style: 2026-10
% course-panel:T11-P01:begin
\par\begingroup
\fontsize{10.5}{13.2}\selectfont
\setlength{\LTleft}{0pt}\setlength{\LTright}{\fill}
\noindent
\begin{tabular}{@{}>{\raggedright\arraybackslash}p{39.459mm}>{\raggedright\arraybackslash}p{39.459mm}>{\raggedright\arraybackslash}p{39.459mm}@{}}
\toprule
\textbf{Person and number} & \textbf{Present indicative} & \textbf{Preterite indicative} \\
\midrule
1 sg. & gemanigfielde & gemanigfielde \\
2 sg. & gemanigfieltst & gemanigfieldest \\
3 sg. & gemanigfielt & gemanigfielde \\
1–3 pl. & gemanigfieldað & gemanigfieldon \\
\bottomrule
\end{tabular}
\par\endgroup
% course-panel:T11-P01:end
% mdtable:T11:end
\reserve{5\baselineskip}
% mdparagraph:P050:begin
\textbf{Tōtwǣman.}
% mdparagraph:P050:end

% teaching: NEW-U07-IND-totwaeman
% mdparagraph:P051:begin
The present stem remains tōtwǣm-; the preterite has tōtwǣmd-. The long ǣ remains in both systems.
% mdparagraph:P051:end

% table: NEW-U07-IND-totwaeman
% mdtable:T12:begin
% course-table-style: 2026-10
% course-panel:T12-P01:begin
\par\begingroup
\fontsize{10.5}{13.2}\selectfont
\setlength{\LTleft}{0pt}\setlength{\LTright}{\fill}
\noindent
\begin{tabular}{@{}>{\raggedright\arraybackslash}p{39.459mm}>{\raggedright\arraybackslash}p{39.459mm}>{\raggedright\arraybackslash}p{39.459mm}@{}}
\toprule
\textbf{Person and number} & \textbf{Present indicative} & \textbf{Preterite indicative} \\
\midrule
1 sg. & tōtwǣme & tōtwǣmde \\
2 sg. & tōtwǣmst & tōtwǣmdest \\
3 sg. & tōtwǣmð & tōtwǣmde \\
1–3 pl. & tōtwǣmað & tōtwǣmdon \\
\bottomrule
\end{tabular}
\par\endgroup
% course-panel:T12-P01:end
% mdtable:T12:end
\reserve{7\baselineskip}
% mdparagraph:P052:begin
\subsection*{Their subjunctives}
% mdparagraph:P052:end

% teaching: NEW-U07-SUBJ-gieman
% teaching: NEW-U07-SUBJ-ricsian
% teaching: NEW-U07-SUBJ-totwaeman
% mdparagraph:P053:begin
Wyrcean’s full subjunctive paradigm appears with weak class I in §2. The remaining new verbs follow in these tables. As before, each shared singular or plural row includes all three persons.
% mdparagraph:P053:end

% table: U07-DISPLAY-SUBJ-6
% mdtable:T13:begin
% course-table-style: 2026-10
% course-panel:T13-P01:begin
\par\begingroup
\fontsize{10.5}{13.2}\selectfont
\setlength{\LTleft}{0pt}\setlength{\LTright}{\fill}
\noindent
\begin{tabular}{@{}>{\raggedright\arraybackslash}p{20.56mm}>{\raggedright\arraybackslash}p{20.748mm}>{\raggedright\arraybackslash}p{24.503mm}>{\raggedright\arraybackslash}p{21.593mm}>{\raggedright\arraybackslash}p{25.348mm}@{}}
\toprule
\textbf{Person and number} & \textbf{gīeman · Present} & \textbf{gīeman · Preterite} & \textbf{rīcsian · Present} & \textbf{rīcsian · Preterite} \\
\midrule
1–3 sg. & gīeme & gīemde & rīcsie & rīcsode \\
1–3 pl. & gīemen & gīemden & rīcsien & rīcsoden \\
\bottomrule
\end{tabular}
\par\endgroup
% course-panel:T13-P01:end
\par\addvspace{5pt}
\par\addvspace{3pt}
% course-caption:T13:1:begin
\textbf{tōtwǣman}\par\nobreak
% course-caption:T13:1:end
% course-panel:T13-P02:begin
\par\begingroup
\fontsize{10.5}{13.2}\selectfont
\setlength{\LTleft}{0pt}\setlength{\LTright}{\fill}
\noindent
\begin{tabular}{@{}>{\raggedright\arraybackslash}p{39.459mm}>{\raggedright\arraybackslash}p{39.459mm}>{\raggedright\arraybackslash}p{39.459mm}@{}}
\toprule
\textbf{Person and number} & \textbf{Present} & \textbf{Preterite} \\
\midrule
1–3 sg. & tōtwǣme & tōtwǣmde \\
1–3 pl. & tōtwǣmen & tōtwǣmden \\
\bottomrule
\end{tabular}
\par\endgroup
% course-panel:T13-P02:end
% mdtable:T13:end
% teaching: NEW-U07-SUBJ-gemanigfieldan
% mdparagraph:P054:begin
Gemanigfieldan has identical present and preterite subjunctive spellings. The table is complete in both tenses; the identity of the letters does not erase the distinction.
% mdparagraph:P054:end

% table: U07-DISPLAY-GEMANIGFIELDAN-SUBJ
% mdtable:T14:begin
% course-table-style: 2026-10
% course-panel:T14-P01:begin
\par\begingroup
\fontsize{10.5}{13.2}\selectfont
\setlength{\LTleft}{0pt}\setlength{\LTright}{\fill}
\noindent
\begin{tabular}{@{}>{\raggedright\arraybackslash}p{39.459mm}>{\raggedright\arraybackslash}p{39.459mm}>{\raggedright\arraybackslash}p{39.459mm}@{}}
\toprule
\textbf{Person and number} & \textbf{Present} & \textbf{Preterite} \\
\midrule
1–3 sg. & gemanigfielde & gemanigfielde \\
1–3 pl. & gemanigfielden & gemanigfielden \\
\bottomrule
\end{tabular}
\par\endgroup
% course-panel:T14-P01:end
% mdtable:T14:end
\reserve{7\baselineskip}
% mdparagraph:P055:begin
\subsection*{The imperative of tōtwǣman}
% mdparagraph:P055:end

% teaching: NEW-U07-IMP-totwaeman
% teaching: NEW-U07-GRAM-U07-TOTWAEMAN-IMPERATIVE
% mdparagraph:P056:begin
Use the singular to address one person and the plural to address several, including two. Preserve the accusative thing divided and fram with the dative thing separated from.
% mdparagraph:P056:end

% table: NEW-U07-IMP-totwaeman
% mdtable:T15:begin
% course-table-style: 2026-10
% course-panel:T15-P01:begin
\par\begingroup
\fontsize{10.5}{13.2}\selectfont
\setlength{\LTleft}{0pt}\setlength{\LTright}{\fill}
\noindent
\begin{tabular}{@{}>{\raggedright\arraybackslash}p{63.342mm}>{\raggedright\arraybackslash}p{57.846mm}@{}}
\toprule
\textbf{Person and number} & \textbf{Command} \\
\midrule
2 sg. & tōtwǣm \\
2 pl. & tōtwǣmað \\
\bottomrule
\end{tabular}
\par\endgroup
% course-panel:T15-P01:end
% mdtable:T15:end
\reserve{7\baselineskip}
% mdparagraph:P057:begin
\subsection*{Practice A. Forms in use}
% mdparagraph:P057:end

% mdparagraph:P058:begin
Use the selected systems. The four entries in A1–A5 must each be labelled by tense and number; the singular or plural represents all three persons.
% mdparagraph:P058:end

\reserve{5\baselineskip}
\par\begin{minipage}{\linewidth}\raggedright
% mdparagraph:P059:begin
\textbf{A1.}
% mdparagraph:P059:end

% item: NEW-U07-A01
% mdparagraph:P060:begin
Give the present singular, present plural, preterite singular, and preterite plural subjunctive. Each singular/\allowbreak{}plural entry represents all three persons; label every entry. Verbs: hīeran, sęllan.
% mdparagraph:P060:end

\end{minipage}\par\addvspace{4pt}

\par\begin{minipage}{\linewidth}\raggedright
% mdparagraph:P061:begin
\textbf{A2.}
% mdparagraph:P061:end

% item: NEW-U07-A02
% mdparagraph:P062:begin
Give the present singular, present plural, preterite singular, and preterite plural subjunctive. Each singular/\allowbreak{}plural entry represents all three persons; label every entry. Verbs: lufian, libban.
% mdparagraph:P062:end

\end{minipage}\par\addvspace{4pt}

\par\begin{minipage}{\linewidth}\raggedright
% mdparagraph:P063:begin
\textbf{A3.}
% mdparagraph:P063:end

% item: NEW-U07-A03
% mdparagraph:P064:begin
Give the present singular, present plural, preterite singular, and preterite plural subjunctive. Each singular/\allowbreak{}plural entry represents all three persons; label every entry. Verbs: bindan, weaxan.
% mdparagraph:P064:end

\end{minipage}\par\addvspace{4pt}

\par\begin{minipage}{\linewidth}\raggedright
% mdparagraph:P065:begin
\textbf{A4.}
% mdparagraph:P065:end

% item: NEW-U07-A04
% mdparagraph:P066:begin
Give the present singular, present plural, preterite singular, and preterite plural subjunctive. Each singular/\allowbreak{}plural entry represents all three persons; label every entry. Verbs: sēon, biddan.
% mdparagraph:P066:end

\end{minipage}\par\addvspace{4pt}

\par\begin{minipage}{\linewidth}\raggedright
% mdparagraph:P067:begin
\textbf{A5.}
% mdparagraph:P067:end

% item: NEW-U07-A05
% mdparagraph:P068:begin
Give the present singular, present plural, preterite singular, and preterite plural subjunctive. Each singular/\allowbreak{}plural entry represents all three persons; label every entry. Verbs: witan, willan.
% mdparagraph:P068:end

\end{minipage}\par\addvspace{4pt}

\par\begin{minipage}{\linewidth}\raggedright
% mdparagraph:P069:begin
\textbf{A6.}
% mdparagraph:P069:end

% item: NEW-U07-A06
% mdparagraph:P070:begin
Give both selected present subjunctive series of wesan/\allowbreak{}bēon and the shared preterite subjunctive, singular and plural. State which person distinctions the forms express.
% mdparagraph:P070:end

\end{minipage}\par\addvspace{4pt}

\par\begin{minipage}{\linewidth}\raggedright
% mdparagraph:P071:begin
\textbf{A7.}
% mdparagraph:P071:end

% item: NEW-U07-A07
% mdparagraph:P072:begin
For gīeman, rīcsian, gemanigfieldan, tōtwǣman, and wyrcean, give the present indicative third singular and preterite indicative third plural. Then give the present subjunctive plural of each.
% mdparagraph:P072:end

\end{minipage}\par\addvspace{4pt}

\par\begin{minipage}{\linewidth}\raggedright
% mdparagraph:P073:begin
\textbf{A8.}
% mdparagraph:P073:end

% item: NEW-U07-A08
% mdparagraph:P074:begin
Contrast hīe bundon with hīe bunden, and hīe lufodon with hīe lufoden. Parse each finite form and explain what the ending contributes. Do not invent a full English sentence from an isolated mood form.
% mdparagraph:P074:end

\end{minipage}\par\addvspace{4pt}

\reserve{7\baselineskip}
% mdparagraph:P075:begin
\section*{4. Basic conditions with gif}
% mdparagraph:P075:end

% teaching: NEW-U07-VOC-gif
% teaching: NEW-U07-VOC-gif-if
% teaching: NEW-U07-GRAM-U07-BASIC-CONDITIONALS
% mdparagraph:P076:begin
\textbf{Gif} means “if”. It introduces a condition; it does not automatically require the subjunctive. An open condition can be stated in the indicative.
% mdparagraph:P076:end

% example: NEW-U07-EXAMPLE-04
% mdparagraph:P077:begin
Gif hē hæfð męte, hē sęlð ðǣm fuglum męte.
% mdparagraph:P077:end

% mdparagraph:P078:begin
If he has food, he gives the birds food.
% mdparagraph:P078:end

% mdparagraph:P079:begin
This is an open condition: having food is presented as a real possibility, without a specially hypothetical setting. Hæfð and sęlð are indicative. The birds are the dative recipients.
% mdparagraph:P079:end

% mdparagraph:P080:begin
A speaker can instead use a present subjunctive in a contingent or tentative condition. The following clause retains the form appropriate to its own meaning.
% mdparagraph:P080:end

% example: NEW-U07-EXAMPLE-05
% mdparagraph:P081:begin
Gif ðū mē helpe, wē magon ðā dēor bindan.
% mdparagraph:P081:end

% mdparagraph:P082:begin
If you help me, we can bind the animals.
% mdparagraph:P082:end

% mdparagraph:P083:begin
Here the speaker tentatively seeks help. Helpe is second-singular present subjunctive; magon is first-plural present indicative, asserting capability on that condition. Mē remains dative after helpan. The main clause need not become subjunctive simply because the condition is subjunctive.
% mdparagraph:P083:end

% teaching: NEW-U07-GRAM-U07-NONPAST-PRETERITE-SUBJUNCTIVE
% mdparagraph:P084:begin
\textbf{Recognition: a preterite form with nonpast reference.} The following setting and interpretation are supplied: he is not coming, and the situation is explicitly unreal and nonpast.
% mdparagraph:P084:end

% example: NEW-U07-EXAMPLE-06
% mdparagraph:P085:begin
Gif hē cōme, ic him hulpe.
% mdparagraph:P085:end

% mdparagraph:P086:begin
If he came, I would help him.
% mdparagraph:P086:end

% mdparagraph:P087:begin
Cōme and hulpe are preterite subjunctives, but this context gives them nonpast reference. The morphological label “preterite” does not itself force past time. Recognize this use in the supplied example; a general system of unreal conditions is not required for independent composition here.
% mdparagraph:P087:end

\reserve{7\baselineskip}
% mdparagraph:P088:begin
\subsection*{Practice M. Mood and meaning}
% mdparagraph:P088:end

% mdparagraph:P089:begin
Distinguish what the form establishes from what the stated context establishes. Where two moods share a spelling, identify the intended reading without declaring the spelling unique.
% mdparagraph:P089:end

\reserve{5\baselineskip}
\par\begin{minipage}{\linewidth}\raggedright
% mdparagraph:P090:begin
\textbf{M1.}
% mdparagraph:P090:end

% item: NEW-U07-M01
% mdparagraph:P091:begin
Translate and contrast Hīe swimmað on ðǣm sǣ and Swimmen hīe on ðǣm sǣ. State what is asserted or requested.
% mdparagraph:P091:end

\end{minipage}\par\addvspace{4pt}

\par\begin{minipage}{\linewidth}\raggedright
% mdparagraph:P092:begin
\textbf{M2.}
% mdparagraph:P092:end

% item: NEW-U07-M02
% mdparagraph:P093:begin
Translate Gāð and Gān wē. Identify the persons addressed or included, and explain why they are not the same command form.
% mdparagraph:P093:end

\end{minipage}\par\addvspace{4pt}

\par\begin{minipage}{\linewidth}\raggedright
% mdparagraph:P094:begin
\textbf{M3.}
% mdparagraph:P094:end

% item: NEW-U07-M03
% mdparagraph:P095:begin
Translate Hē mæg swimman, Hē mōt swimman, and Swimme hē. Distinguish capability, permission, and the speaker’s wish.
% mdparagraph:P095:end

\end{minipage}\par\addvspace{4pt}

\par\begin{minipage}{\linewidth}\raggedright
% mdparagraph:P096:begin
\textbf{M4.}
% mdparagraph:P096:end

% item: NEW-U07-M04
% mdparagraph:P097:begin
An open condition describes a real possibility without presenting it as specially hypothetical: Gif hē hæfð męte, hē sęlð ðǣm fuglum męte. Parse and translate the clause pair. Does gif itself require subjunctive morphology?
% mdparagraph:P097:end

\end{minipage}\par\addvspace{4pt}

\par\begin{minipage}{\linewidth}\raggedright
% mdparagraph:P098:begin
\textbf{M5.}
% mdparagraph:P098:end

% item: NEW-U07-M05
% mdparagraph:P099:begin
A speaker tentatively seeks help: Gif ðū mē helpe, wē magon ðā dēor bindan. Parse helpe and magon, translate the sentence, and explain the difference between their roles.
% mdparagraph:P099:end

\end{minipage}\par\addvspace{4pt}

\par\begin{minipage}{\linewidth}\raggedright
% mdparagraph:P100:begin
\textbf{M6.}
% mdparagraph:P100:end

% item: NEW-U07-M06
% mdparagraph:P101:begin
The setting is explicitly unreal and nonpast: he is not coming. The supplied interpretation of Gif hē cōme, ic him hulpe is “If he came, I would help him.” Identify the tense and mood of cōme and hulpe. Explain why their morphological label alone does not force past time.
% mdparagraph:P101:end

\end{minipage}\par\addvspace{4pt}

\par\begin{minipage}{\linewidth}\raggedright
% mdparagraph:P102:begin
\textbf{M7.}
% mdparagraph:P102:end

% item: NEW-U07-M07
% mdparagraph:P103:begin
The narrator asserts a past wish in Hē wolde faran. Parse wolde according to that intended assertion and translate it. State the other mood whose form has the same spelling, and explain why the spelling alone does not require the English word “would.”
% mdparagraph:P103:end

\end{minipage}\par\addvspace{4pt}

\par\begin{minipage}{\linewidth}\raggedright
% mdparagraph:P104:begin
\textbf{M8.}
% mdparagraph:P104:end

% item: NEW-U07-M08
% mdparagraph:P105:begin
Give a third-person singular wish corresponding to Hē leofað, and a first-person plural exhortation corresponding to Wē helpað ðǣm męnn. Preserve the lexical meaning and government.
% mdparagraph:P105:end

\end{minipage}\par\addvspace{4pt}

\reserve{7\baselineskip}
% mdparagraph:P106:begin
\section*{5. Participles}
% mdparagraph:P106:end

% teaching: NEW-U07-GRAM-U07-PARTICIPLE-FORMATION
% mdparagraph:P107:begin
A \textbf{participle} is a nonfinite verb form. The present participle in \textbf{-ende} describes an activity; a past participle can describe a resulting condition or enter a passive construction. Neither participle by itself is a finite present or preterite verb.
% mdparagraph:P107:end

% mdparagraph:P108:begin
For the selected present participles, retain the present stem before -ende. Verbs in -ian retain i: nęrian gives nęriende and lufian gives lufiende. Gefręmman retains double mm in its present participle. The strong past participle uses its own principal-part stem and -en, as in bunden, boren, and cweden; it is not made by adding a weak dental ending to the infinitive. The weak past participles have their selected dental forms, including hīered, lufod, and seald.
% mdparagraph:P108:end

% mdparagraph:P109:begin
The following twelve pairs are the formation assignment. Preserve the actual lexical prefix: geblētsod and gehālgod belong to their prefixed verbs; bunden and hāten are the selected unprefixed forms. A dictionary’s optional prefix notation is not a word to copy into a sentence, and it is not a rule for adding ge- to arbitrary forms.
% mdparagraph:P109:end

% teaching: NEW-U07-PART-bindan
% teaching: NEW-U07-PART-hieran
% teaching: NEW-U07-PART-lufian
% teaching: NEW-U07-PART-sellan
% teaching: NEW-U07-PART-hatan
% teaching: NEW-U07-PART-gemanigfieldan
% teaching: NEW-U07-PART-nerian
% teaching: NEW-U07-PART-cwethan
% teaching: NEW-U07-PART-beran
% teaching: NEW-U07-PART-gehalgian
% teaching: NEW-U07-PART-gebletsian
% teaching: NEW-U07-PART-gefremman
% table: U07-DISPLAY-PART-PAIRS
% mdtable:T16:begin
\par\begingroup
\fontsize{10.5}{13.2}\selectfont
\setlength{\LTleft}{0pt}\setlength{\LTright}{\fill}
\noindent
\begin{tabular}{@{}>{\raggedright\arraybackslash}p{37.45mm}>{\raggedright\arraybackslash}p{45.21mm}>{\raggedright\arraybackslash}p{35.71mm}@{}}
\toprule
\textbf{Lemma} & \textbf{Present participle} & \textbf{Past participle} \\
\midrule
bindan & bindende & bunden \\
hīeran & hīerende & hīered \\
lufian & lufiende & lufod \\
sęllan & sęllende & seald \\
hātan & hātende & hāten \\
gemanigfieldan & gemanigfieldende & gemanigfield \\
nęrian & nęriende & nęred \\
cweðan & cweðende & cweden \\
beran & berende & boren \\
gehālgian & gehālgiende & gehālgod \\
geblētsian & geblētsiende & geblētsod \\
gefręmman & gefręmmende & gefręmed \\
\bottomrule
\end{tabular}
\par\endgroup

% mdtable:T16:end
% mdparagraph:P110:begin
The weak preterite is not a universal past-participle recipe. Hīerde and hīered differ in the placement of e; lufode corresponds to lufod, and sealde to seald. Gefręmed has single m, unlike the present gefręmmende. Learn the actual principal part with its class rather than merely removing the final letter of every past form.
% mdparagraph:P110:end

% teaching: NEW-U07-GRAM-U07-PARTICIPLE-AGREEMENT
\reserve{7\baselineskip}
% mdparagraph:P111:begin
\subsection*{Participles used as adjectives}
% mdparagraph:P111:end

% mdparagraph:P112:begin
An attributive participle agrees with its noun in \textbf{gender, number, and case}. Use the strong or weak declension under the conditions already learned for adjectives: a definite article calls for the weak attributive form; an otherwise undetermined adjective ordinarily takes the strong form. No new instrumental use is added here.
% mdparagraph:P112:end

% mdparagraph:P113:begin
The five complete models below give all four cases in all three genders and both numbers, in both strong and weak inflection. The present participle follows the grēne-type adjective pattern. It is not an agent noun in -end.
% mdparagraph:P113:end

% example: NEW-U07-EXAMPLE-07
% mdparagraph:P114:begin
Se bindenda mann is gōd.
% mdparagraph:P114:end

% mdparagraph:P115:begin
The binding man is good.
% mdparagraph:P115:end

% mdparagraph:P116:begin
Bindenda is weak masculine nominative singular after se, modifying mann. Gōd is the separate predicate adjective.
% mdparagraph:P116:end

\reserve{7\baselineskip}
% mdparagraph:P117:begin
\subsection*{Bindende: present participle}
% mdparagraph:P117:end

% teaching: NEW-U07-PART-ENDE
% mdparagraph:P118:begin
Retain the participial stem bindend- before consonantal endings, and observe the selected -e/\allowbreak{}-u distribution in the strong forms. The weak nominative singular has masculine -a beside feminine and neuter -e.
% mdparagraph:P118:end

\reserve{9\baselineskip}
% mdparagraph:P119:begin
\textbf{Strong singular.}
% mdparagraph:P119:end

% table: NEW-U07-PART-ENDE-1
% mdtable:T17:begin
% course-table-style: 2026-10
% course-panel:T17-P01:begin
\par\begingroup
\fontsize{10.5}{13.2}\selectfont
\setlength{\LTleft}{0pt}\setlength{\LTright}{\fill}
\noindent
\begin{tabular}{@{}>{\raggedright\arraybackslash}p{22.0mm}>{\raggedright\arraybackslash}p{31.188mm}>{\raggedright\arraybackslash}p{31.188mm}>{\raggedright\arraybackslash}p{31.188mm}@{}}
\toprule
\textbf{Case} & \textbf{Masculine} & \textbf{Neuter} & \textbf{Feminine} \\
\midrule
Nominative & bindende & bindende & bindendu \\
Genitive & bindendes & bindendes & bindendre \\
Dative & bindendum & bindendum & bindendre \\
Accusative & bindendne & bindende & bindende \\
\bottomrule
\end{tabular}
\par\endgroup
% course-panel:T17-P01:end
% mdtable:T17:end
\reserve{9\baselineskip}
% mdparagraph:P120:begin
\textbf{Strong plural.}
% mdparagraph:P120:end

% table: NEW-U07-PART-ENDE-2
% mdtable:T18:begin
% course-table-style: 2026-10
% course-panel:T18-P01:begin
\par\begingroup
\fontsize{10.5}{13.2}\selectfont
\setlength{\LTleft}{0pt}\setlength{\LTright}{\fill}
\noindent
\begin{tabular}{@{}>{\raggedright\arraybackslash}p{22.0mm}>{\raggedright\arraybackslash}p{31.188mm}>{\raggedright\arraybackslash}p{31.188mm}>{\raggedright\arraybackslash}p{31.188mm}@{}}
\toprule
\textbf{Case} & \textbf{Masculine} & \textbf{Neuter} & \textbf{Feminine} \\
\midrule
Nominative & bindende & bindendu & bindenda \\
Genitive & bindendra & bindendra & bindendra \\
Dative & bindendum & bindendum & bindendum \\
Accusative & bindende & bindendu & bindenda \\
\bottomrule
\end{tabular}
\par\endgroup
% course-panel:T18-P01:end
% mdtable:T18:end
\reserve{9\baselineskip}
% mdparagraph:P121:begin
\textbf{Weak singular.}
% mdparagraph:P121:end

% table: NEW-U07-PART-ENDE-3
% mdtable:T19:begin
% course-table-style: 2026-10
% course-panel:T19-P01:begin
\par\begingroup
\fontsize{10.5}{13.2}\selectfont
\setlength{\LTleft}{0pt}\setlength{\LTright}{\fill}
\noindent
\begin{tabular}{@{}>{\raggedright\arraybackslash}p{22.0mm}>{\raggedright\arraybackslash}p{31.188mm}>{\raggedright\arraybackslash}p{31.188mm}>{\raggedright\arraybackslash}p{31.188mm}@{}}
\toprule
\textbf{Case} & \textbf{Masculine} & \textbf{Neuter} & \textbf{Feminine} \\
\midrule
Nominative & bindenda & bindende & bindende \\
Genitive & bindendan & bindendan & bindendan \\
Dative & bindendan & bindendan & bindendan \\
Accusative & bindendan & bindende & bindendan \\
\bottomrule
\end{tabular}
\par\endgroup
% course-panel:T19-P01:end
% mdtable:T19:end
\reserve{9\baselineskip}
% mdparagraph:P122:begin
\textbf{Weak plural.}
% mdparagraph:P122:end

% table: NEW-U07-PART-ENDE-4
% mdtable:T20:begin
% course-table-style: 2026-10
% course-panel:T20-P01:begin
\par\begingroup
\fontsize{10.5}{13.2}\selectfont
\setlength{\LTleft}{0pt}\setlength{\LTright}{\fill}
\noindent
\begin{tabular}{@{}>{\raggedright\arraybackslash}p{22.0mm}>{\raggedright\arraybackslash}p{31.188mm}>{\raggedright\arraybackslash}p{31.188mm}>{\raggedright\arraybackslash}p{31.188mm}@{}}
\toprule
\textbf{Case} & \textbf{Masculine} & \textbf{Neuter} & \textbf{Feminine} \\
\midrule
Nominative & bindendan & bindendan & bindendan \\
Genitive & bindendra & bindendra & bindendra \\
Dative & bindendum & bindendum & bindendum \\
Accusative & bindendan & bindendan & bindendan \\
\bottomrule
\end{tabular}
\par\endgroup
% course-panel:T20-P01:end
% mdtable:T20:end
% mdparagraph:P123:begin
\textbf{Variants.} In the strong feminine plural nominative/\allowbreak{}accusative, bindende is also found beside selected bindenda; in the corresponding neuter plural cells, bindende is an alternative to bindendu. Keep the selected forms when a task asks for the model. These variations do not make the participle an uninflected word in every attributive use.
% mdparagraph:P123:end

% teaching: NEW-U07-GRAM-U07-PARTICIPLE-SYNCOPE
\reserve{7\baselineskip}
% mdparagraph:P124:begin
\subsection*{Bunden: strong past participle}
% mdparagraph:P124:end

% teaching: NEW-U07-PART-BUNDEN
% mdparagraph:P125:begin
The table keeps the full stem bunden- so that the endings remain visible. Its unstressed e can disappear before a vowel-initial ending: this is syncope. Preserve the full stem before -ne, -re, and -ra; bundenne therefore has double n.
% mdparagraph:P125:end

\reserve{9\baselineskip}
% mdparagraph:P126:begin
\textbf{Strong singular.}
% mdparagraph:P126:end

% table: NEW-U07-PART-BUNDEN-1
% mdtable:T21:begin
% course-table-style: 2026-10
% course-panel:T21-P01:begin
\par\begingroup
\fontsize{10.5}{13.2}\selectfont
\setlength{\LTleft}{0pt}\setlength{\LTright}{\fill}
\noindent
\begin{tabular}{@{}>{\raggedright\arraybackslash}p{22.0mm}>{\raggedright\arraybackslash}p{31.188mm}>{\raggedright\arraybackslash}p{31.188mm}>{\raggedright\arraybackslash}p{31.188mm}@{}}
\toprule
\textbf{Case} & \textbf{Masculine} & \textbf{Neuter} & \textbf{Feminine} \\
\midrule
Nominative & bunden & bunden & bundenu \\
Genitive & bundenes & bundenes & bundenre \\
Dative & bundenum & bundenum & bundenre \\
Accusative & bundenne & bunden & bundene \\
\bottomrule
\end{tabular}
\par\endgroup
% course-panel:T21-P01:end
% mdtable:T21:end
\reserve{9\baselineskip}
% mdparagraph:P127:begin
\textbf{Strong plural.}
% mdparagraph:P127:end

% table: NEW-U07-PART-BUNDEN-2
% mdtable:T22:begin
% course-table-style: 2026-10
% course-panel:T22-P01:begin
\par\begingroup
\fontsize{10.5}{13.2}\selectfont
\setlength{\LTleft}{0pt}\setlength{\LTright}{\fill}
\noindent
\begin{tabular}{@{}>{\raggedright\arraybackslash}p{22.0mm}>{\raggedright\arraybackslash}p{31.188mm}>{\raggedright\arraybackslash}p{31.188mm}>{\raggedright\arraybackslash}p{31.188mm}@{}}
\toprule
\textbf{Case} & \textbf{Masculine} & \textbf{Neuter} & \textbf{Feminine} \\
\midrule
Nominative & bundene & bundenu & bundena \\
Genitive & bundenra & bundenra & bundenra \\
Dative & bundenum & bundenum & bundenum \\
Accusative & bundene & bundenu & bundena \\
\bottomrule
\end{tabular}
\par\endgroup
% course-panel:T22-P01:end
% mdtable:T22:end
\reserve{9\baselineskip}
% mdparagraph:P128:begin
\textbf{Weak singular.}
% mdparagraph:P128:end

% table: NEW-U07-PART-BUNDEN-3
% mdtable:T23:begin
% course-table-style: 2026-10
% course-panel:T23-P01:begin
\par\begingroup
\fontsize{10.5}{13.2}\selectfont
\setlength{\LTleft}{0pt}\setlength{\LTright}{\fill}
\noindent
\begin{tabular}{@{}>{\raggedright\arraybackslash}p{22.0mm}>{\raggedright\arraybackslash}p{31.188mm}>{\raggedright\arraybackslash}p{31.188mm}>{\raggedright\arraybackslash}p{31.188mm}@{}}
\toprule
\textbf{Case} & \textbf{Masculine} & \textbf{Neuter} & \textbf{Feminine} \\
\midrule
Nominative & bundena & bundene & bundene \\
Genitive & bundenan & bundenan & bundenan \\
Dative & bundenan & bundenan & bundenan \\
Accusative & bundenan & bundene & bundenan \\
\bottomrule
\end{tabular}
\par\endgroup
% course-panel:T23-P01:end
% mdtable:T23:end
\reserve{9\baselineskip}
% mdparagraph:P129:begin
\textbf{Weak plural.}
% mdparagraph:P129:end

% table: NEW-U07-PART-BUNDEN-4
% mdtable:T24:begin
% course-table-style: 2026-10
% course-panel:T24-P01:begin
\par\begingroup
\fontsize{10.5}{13.2}\selectfont
\setlength{\LTleft}{0pt}\setlength{\LTright}{\fill}
\noindent
\begin{tabular}{@{}>{\raggedright\arraybackslash}p{22.0mm}>{\raggedright\arraybackslash}p{31.188mm}>{\raggedright\arraybackslash}p{31.188mm}>{\raggedright\arraybackslash}p{31.188mm}@{}}
\toprule
\textbf{Case} & \textbf{Masculine} & \textbf{Neuter} & \textbf{Feminine} \\
\midrule
Nominative & bundenan & bundenan & bundenan \\
Genitive & bundenra & bundenra & bundenra \\
Dative & bundenum & bundenum & bundenum \\
Accusative & bundenan & bundenan & bundenan \\
\bottomrule
\end{tabular}
\par\endgroup
% course-panel:T24-P01:end
% mdtable:T24:end
% mdparagraph:P130:begin
\textbf{Syncope and variants.} Bundenes/\allowbreak{}bundnes, bundenum/\allowbreak{}bundnum, bundene/\allowbreak{}bundne, bundena/\allowbreak{}bundna, and bundenan/\allowbreak{}bundnan are the corresponding full and syncopated forms. Apply each pair only to the cells with that full form. Strong feminine plural nominative/\allowbreak{}accusative also permits bundene beside bundena. Strong feminine nominative singular permits bunden beside bundenu; strong neuter plural nominative/\allowbreak{}accusative permits bunden or bundene beside bundenu. The selected -u forms complete the model; they are not claims that every such cell is frequent in continuous texts.
% mdparagraph:P130:end

\reserve{7\baselineskip}
% mdparagraph:P131:begin
\subsection*{Hīered: weak past participle with syncope}
% mdparagraph:P131:end

% teaching: NEW-U07-PART-HIERED
% mdparagraph:P132:begin
The full stem hīered- makes the relation to the dental participle clear. Keep it before -ne, -re, and -ra. Before vowel-initial endings its unstressed e can disappear, producing hīerd-.
% mdparagraph:P132:end

\reserve{9\baselineskip}
% mdparagraph:P133:begin
\textbf{Strong singular.}
% mdparagraph:P133:end

% table: NEW-U07-PART-HIERED-1
% mdtable:T25:begin
% course-table-style: 2026-10
% course-panel:T25-P01:begin
\par\begingroup
\fontsize{10.5}{13.2}\selectfont
\setlength{\LTleft}{0pt}\setlength{\LTright}{\fill}
\noindent
\begin{tabular}{@{}>{\raggedright\arraybackslash}p{22.0mm}>{\raggedright\arraybackslash}p{31.188mm}>{\raggedright\arraybackslash}p{31.188mm}>{\raggedright\arraybackslash}p{31.188mm}@{}}
\toprule
\textbf{Case} & \textbf{Masculine} & \textbf{Neuter} & \textbf{Feminine} \\
\midrule
Nominative & hīered & hīered & hīeredu \\
Genitive & hīeredes & hīeredes & hīeredre \\
Dative & hīeredum & hīeredum & hīeredre \\
Accusative & hīeredne & hīered & hīerede \\
\bottomrule
\end{tabular}
\par\endgroup
% course-panel:T25-P01:end
% mdtable:T25:end
\reserve{9\baselineskip}
% mdparagraph:P134:begin
\textbf{Strong plural.}
% mdparagraph:P134:end

% table: NEW-U07-PART-HIERED-2
% mdtable:T26:begin
% course-table-style: 2026-10
% course-panel:T26-P01:begin
\par\begingroup
\fontsize{10.5}{13.2}\selectfont
\setlength{\LTleft}{0pt}\setlength{\LTright}{\fill}
\noindent
\begin{tabular}{@{}>{\raggedright\arraybackslash}p{22.0mm}>{\raggedright\arraybackslash}p{31.188mm}>{\raggedright\arraybackslash}p{31.188mm}>{\raggedright\arraybackslash}p{31.188mm}@{}}
\toprule
\textbf{Case} & \textbf{Masculine} & \textbf{Neuter} & \textbf{Feminine} \\
\midrule
Nominative & hīerede & hīeredu & hīereda \\
Genitive & hīeredra & hīeredra & hīeredra \\
Dative & hīeredum & hīeredum & hīeredum \\
Accusative & hīerede & hīeredu & hīereda \\
\bottomrule
\end{tabular}
\par\endgroup
% course-panel:T26-P01:end
% mdtable:T26:end
\reserve{9\baselineskip}
% mdparagraph:P135:begin
\textbf{Weak singular.}
% mdparagraph:P135:end

% table: NEW-U07-PART-HIERED-3
% mdtable:T27:begin
% course-table-style: 2026-10
% course-panel:T27-P01:begin
\par\begingroup
\fontsize{10.5}{13.2}\selectfont
\setlength{\LTleft}{0pt}\setlength{\LTright}{\fill}
\noindent
\begin{tabular}{@{}>{\raggedright\arraybackslash}p{22.0mm}>{\raggedright\arraybackslash}p{31.188mm}>{\raggedright\arraybackslash}p{31.188mm}>{\raggedright\arraybackslash}p{31.188mm}@{}}
\toprule
\textbf{Case} & \textbf{Masculine} & \textbf{Neuter} & \textbf{Feminine} \\
\midrule
Nominative & hīereda & hīerede & hīerede \\
Genitive & hīeredan & hīeredan & hīeredan \\
Dative & hīeredan & hīeredan & hīeredan \\
Accusative & hīeredan & hīerede & hīeredan \\
\bottomrule
\end{tabular}
\par\endgroup
% course-panel:T27-P01:end
% mdtable:T27:end
\reserve{9\baselineskip}
% mdparagraph:P136:begin
\textbf{Weak plural.}
% mdparagraph:P136:end

% table: NEW-U07-PART-HIERED-4
% mdtable:T28:begin
% course-table-style: 2026-10
% course-panel:T28-P01:begin
\par\begingroup
\fontsize{10.5}{13.2}\selectfont
\setlength{\LTleft}{0pt}\setlength{\LTright}{\fill}
\noindent
\begin{tabular}{@{}>{\raggedright\arraybackslash}p{22.0mm}>{\raggedright\arraybackslash}p{31.188mm}>{\raggedright\arraybackslash}p{31.188mm}>{\raggedright\arraybackslash}p{31.188mm}@{}}
\toprule
\textbf{Case} & \textbf{Masculine} & \textbf{Neuter} & \textbf{Feminine} \\
\midrule
Nominative & hīeredan & hīeredan & hīeredan \\
Genitive & hīeredra & hīeredra & hīeredra \\
Dative & hīeredum & hīeredum & hīeredum \\
Accusative & hīeredan & hīeredan & hīeredan \\
\bottomrule
\end{tabular}
\par\endgroup
% course-panel:T28-P01:end
% mdtable:T28:end
% mdparagraph:P137:begin
\textbf{Syncope and variants.} Hīeredes/\allowbreak{}hīerdes, hīeredum/\allowbreak{}hīerdum, hīerede/\allowbreak{}hīerde, hīereda/\allowbreak{}hīerda, and hīeredan/\allowbreak{}hīerdan give the corresponding alternatives. Strong feminine plural nominative/\allowbreak{}accusative also permits hīerede beside hīereda. Strong feminine nominative singular permits hīered beside hīeredu; strong neuter plural nominative/\allowbreak{}accusative permits hīered or hīerede beside hīeredu. As in the preceding model, the -u cells are regular paradigm completions with their alternatives stated explicitly.
% mdparagraph:P137:end

\reserve{7\baselineskip}
% mdparagraph:P138:begin
\subsection*{Lufod: class II past participle}
% mdparagraph:P138:end

% teaching: NEW-U07-PART-LUFOD
% mdparagraph:P139:begin
Retain o throughout the selected model. Unlike the preceding -u models, selected lufod has no ending in the strong feminine nominative singular or neuter plural nominative/\allowbreak{}accusative. The remaining endings display the familiar strong/\allowbreak{}weak contrasts.
% mdparagraph:P139:end

\reserve{9\baselineskip}
% mdparagraph:P140:begin
\textbf{Strong singular.}
% mdparagraph:P140:end

% table: NEW-U07-PART-LUFOD-1
% mdtable:T29:begin
% course-table-style: 2026-10
% course-panel:T29-P01:begin
\par\begingroup
\fontsize{10.5}{13.2}\selectfont
\setlength{\LTleft}{0pt}\setlength{\LTright}{\fill}
\noindent
\begin{tabular}{@{}>{\raggedright\arraybackslash}p{22.0mm}>{\raggedright\arraybackslash}p{31.188mm}>{\raggedright\arraybackslash}p{31.188mm}>{\raggedright\arraybackslash}p{31.188mm}@{}}
\toprule
\textbf{Case} & \textbf{Masculine} & \textbf{Neuter} & \textbf{Feminine} \\
\midrule
Nominative & lufod & lufod & lufod \\
Genitive & lufodes & lufodes & lufodre \\
Dative & lufodum & lufodum & lufodre \\
Accusative & lufodne & lufod & lufode \\
\bottomrule
\end{tabular}
\par\endgroup
% course-panel:T29-P01:end
% mdtable:T29:end
\reserve{9\baselineskip}
% mdparagraph:P141:begin
\textbf{Strong plural.}
% mdparagraph:P141:end

% table: NEW-U07-PART-LUFOD-2
% mdtable:T30:begin
% course-table-style: 2026-10
% course-panel:T30-P01:begin
\par\begingroup
\fontsize{10.5}{13.2}\selectfont
\setlength{\LTleft}{0pt}\setlength{\LTright}{\fill}
\noindent
\begin{tabular}{@{}>{\raggedright\arraybackslash}p{22.0mm}>{\raggedright\arraybackslash}p{31.188mm}>{\raggedright\arraybackslash}p{31.188mm}>{\raggedright\arraybackslash}p{31.188mm}@{}}
\toprule
\textbf{Case} & \textbf{Masculine} & \textbf{Neuter} & \textbf{Feminine} \\
\midrule
Nominative & lufode & lufod & lufoda \\
Genitive & lufodra & lufodra & lufodra \\
Dative & lufodum & lufodum & lufodum \\
Accusative & lufode & lufod & lufoda \\
\bottomrule
\end{tabular}
\par\endgroup
% course-panel:T30-P01:end
% mdtable:T30:end
\reserve{9\baselineskip}
% mdparagraph:P142:begin
\textbf{Weak singular.}
% mdparagraph:P142:end

% table: NEW-U07-PART-LUFOD-3
% mdtable:T31:begin
% course-table-style: 2026-10
% course-panel:T31-P01:begin
\par\begingroup
\fontsize{10.5}{13.2}\selectfont
\setlength{\LTleft}{0pt}\setlength{\LTright}{\fill}
\noindent
\begin{tabular}{@{}>{\raggedright\arraybackslash}p{22.0mm}>{\raggedright\arraybackslash}p{31.188mm}>{\raggedright\arraybackslash}p{31.188mm}>{\raggedright\arraybackslash}p{31.188mm}@{}}
\toprule
\textbf{Case} & \textbf{Masculine} & \textbf{Neuter} & \textbf{Feminine} \\
\midrule
Nominative & lufoda & lufode & lufode \\
Genitive & lufodan & lufodan & lufodan \\
Dative & lufodan & lufodan & lufodan \\
Accusative & lufodan & lufode & lufodan \\
\bottomrule
\end{tabular}
\par\endgroup
% course-panel:T31-P01:end
% mdtable:T31:end
\reserve{9\baselineskip}
% mdparagraph:P143:begin
\textbf{Weak plural.}
% mdparagraph:P143:end

% table: NEW-U07-PART-LUFOD-4
% mdtable:T32:begin
% course-table-style: 2026-10
% course-panel:T32-P01:begin
\par\begingroup
\fontsize{10.5}{13.2}\selectfont
\setlength{\LTleft}{0pt}\setlength{\LTright}{\fill}
\noindent
\begin{tabular}{@{}>{\raggedright\arraybackslash}p{22.0mm}>{\raggedright\arraybackslash}p{31.188mm}>{\raggedright\arraybackslash}p{31.188mm}>{\raggedright\arraybackslash}p{31.188mm}@{}}
\toprule
\textbf{Case} & \textbf{Masculine} & \textbf{Neuter} & \textbf{Feminine} \\
\midrule
Nominative & lufodan & lufodan & lufodan \\
Genitive & lufodra & lufodra & lufodra \\
Dative & lufodum & lufodum & lufodum \\
Accusative & lufodan & lufodan & lufodan \\
\bottomrule
\end{tabular}
\par\endgroup
% course-panel:T32-P01:end
% mdtable:T32:end
% mdparagraph:P144:begin
\textbf{Variants.} Strong feminine plural nominative/\allowbreak{}accusative lufode is an alternative to lufoda; strong neuter plural nominative/\allowbreak{}accusative lufode is an alternative to lufod. These are particular agreement variants, not a replacement of every final syllable by -e.
% mdparagraph:P144:end

\reserve{7\baselineskip}
% mdparagraph:P145:begin
\subsection*{Seald: monosyllabic dental past participle}
% mdparagraph:P145:end

% teaching: NEW-U07-PART-SEALD
% mdparagraph:P146:begin
Seald has no unstressed stem vowel to remove. Add the printed endings directly to seald-. The strong feminine nominative singular and neuter plural nominative/\allowbreak{}accusative are unextended in the selected model.
% mdparagraph:P146:end

\reserve{9\baselineskip}
% mdparagraph:P147:begin
\textbf{Strong singular.}
% mdparagraph:P147:end

% table: NEW-U07-PART-SEALD-1
% mdtable:T33:begin
% course-table-style: 2026-10
% course-panel:T33-P01:begin
\par\begingroup
\fontsize{10.5}{13.2}\selectfont
\setlength{\LTleft}{0pt}\setlength{\LTright}{\fill}
\noindent
\begin{tabular}{@{}>{\raggedright\arraybackslash}p{22.0mm}>{\raggedright\arraybackslash}p{31.188mm}>{\raggedright\arraybackslash}p{31.188mm}>{\raggedright\arraybackslash}p{31.188mm}@{}}
\toprule
\textbf{Case} & \textbf{Masculine} & \textbf{Neuter} & \textbf{Feminine} \\
\midrule
Nominative & seald & seald & seald \\
Genitive & sealdes & sealdes & sealdre \\
Dative & sealdum & sealdum & sealdre \\
Accusative & sealdne & seald & sealde \\
\bottomrule
\end{tabular}
\par\endgroup
% course-panel:T33-P01:end
% mdtable:T33:end
\reserve{9\baselineskip}
% mdparagraph:P148:begin
\textbf{Strong plural.}
% mdparagraph:P148:end

% table: NEW-U07-PART-SEALD-2
% mdtable:T34:begin
% course-table-style: 2026-10
% course-panel:T34-P01:begin
\par\begingroup
\fontsize{10.5}{13.2}\selectfont
\setlength{\LTleft}{0pt}\setlength{\LTright}{\fill}
\noindent
\begin{tabular}{@{}>{\raggedright\arraybackslash}p{22.0mm}>{\raggedright\arraybackslash}p{31.188mm}>{\raggedright\arraybackslash}p{31.188mm}>{\raggedright\arraybackslash}p{31.188mm}@{}}
\toprule
\textbf{Case} & \textbf{Masculine} & \textbf{Neuter} & \textbf{Feminine} \\
\midrule
Nominative & sealde & seald & sealda \\
Genitive & sealdra & sealdra & sealdra \\
Dative & sealdum & sealdum & sealdum \\
Accusative & sealde & seald & sealda \\
\bottomrule
\end{tabular}
\par\endgroup
% course-panel:T34-P01:end
% mdtable:T34:end
\reserve{9\baselineskip}
% mdparagraph:P149:begin
\textbf{Weak singular.}
% mdparagraph:P149:end

% table: NEW-U07-PART-SEALD-3
% mdtable:T35:begin
% course-table-style: 2026-10
% course-panel:T35-P01:begin
\par\begingroup
\fontsize{10.5}{13.2}\selectfont
\setlength{\LTleft}{0pt}\setlength{\LTright}{\fill}
\noindent
\begin{tabular}{@{}>{\raggedright\arraybackslash}p{22.0mm}>{\raggedright\arraybackslash}p{31.188mm}>{\raggedright\arraybackslash}p{31.188mm}>{\raggedright\arraybackslash}p{31.188mm}@{}}
\toprule
\textbf{Case} & \textbf{Masculine} & \textbf{Neuter} & \textbf{Feminine} \\
\midrule
Nominative & sealda & sealde & sealde \\
Genitive & sealdan & sealdan & sealdan \\
Dative & sealdan & sealdan & sealdan \\
Accusative & sealdan & sealde & sealdan \\
\bottomrule
\end{tabular}
\par\endgroup
% course-panel:T35-P01:end
% mdtable:T35:end
\reserve{9\baselineskip}
% mdparagraph:P150:begin
\textbf{Weak plural.}
% mdparagraph:P150:end

% table: NEW-U07-PART-SEALD-4
% mdtable:T36:begin
% course-table-style: 2026-10
% course-panel:T36-P01:begin
\par\begingroup
\fontsize{10.5}{13.2}\selectfont
\setlength{\LTleft}{0pt}\setlength{\LTright}{\fill}
\noindent
\begin{tabular}{@{}>{\raggedright\arraybackslash}p{22.0mm}>{\raggedright\arraybackslash}p{31.188mm}>{\raggedright\arraybackslash}p{31.188mm}>{\raggedright\arraybackslash}p{31.188mm}@{}}
\toprule
\textbf{Case} & \textbf{Masculine} & \textbf{Neuter} & \textbf{Feminine} \\
\midrule
Nominative & sealdan & sealdan & sealdan \\
Genitive & sealdra & sealdra & sealdra \\
Dative & sealdum & sealdum & sealdum \\
Accusative & sealdan & sealdan & sealdan \\
\bottomrule
\end{tabular}
\par\endgroup
% course-panel:T36-P01:end
% mdtable:T36:end
% mdparagraph:P151:begin
\textbf{Variants.} Strong feminine plural nominative/\allowbreak{}accusative sealde is an alternative to sealda; strong neuter plural nominative/\allowbreak{}accusative sealde is an alternative to seald.
% mdparagraph:P151:end

\reserve{7\baselineskip}
% mdparagraph:P152:begin
\section*{6. Elementary passive constructions}
% mdparagraph:P152:end

% teaching: NEW-U07-VOC-wesan-passive
% teaching: NEW-U07-VOC-weorthan-1
% teaching: NEW-U07-GRAM-U07-PASSIVE
% mdparagraph:P153:begin
The \textbf{passive} makes the affected person or thing the subject. With the transitive verbs used here, an active accusative object corresponds to a nominative passive subject. Use a finite form of \textbf{wesan/\allowbreak{}bēon} or \textbf{weorðan} with the past participle. The auxiliary agrees with the new subject in person and number and carries the tense and mood. The agent can be omitted.
% mdparagraph:P153:end

% example: NEW-U07-EXAMPLE-09
% mdparagraph:P154:begin
Se mann band ðæt dēor. Ðæt dēor wearð bunden.
% mdparagraph:P154:end

% mdparagraph:P155:begin
The man bound the animal. The animal was bound.
% mdparagraph:P155:end

% mdparagraph:P156:begin
Dēor is accusative in the active clause and nominative in the passive, although its neuter singular spelling is unchanged. Wearð is the singular preterite auxiliary; bunden is the past participle.
% mdparagraph:P156:end

% mdparagraph:P157:begin
Wesan/\allowbreak{}bēon often presents a state or result; weorðan often presents an event or becoming. These are useful tendencies, not an absolute division. Do not assign a fixed modern English tense or aspect to every auxiliary–participle combination. The past participle does not carry the finite tense by itself.
% mdparagraph:P157:end

% teaching: NEW-U07-PRED-PART
% teaching: NEW-U07-GRAM-U07-PREDICATE-PARTICIPLE-VARIANTS
\reserve{7\baselineskip}
% mdparagraph:P158:begin
\subsection*{Verbal predicates and adjective agreement}
% mdparagraph:P158:end

% mdparagraph:P159:begin
In a verbal passive the past participle may be \textbf{uninflected}, or it may take \textbf{strong masculine-like endings}. This differs from the agreement rule for an attributive adjective. In particular, a neuter plural subject does not require the attributive neuter-plural -u form in every verbal passive.
% mdparagraph:P159:end

% example: NEW-U07-EXAMPLE-08
% mdparagraph:P160:begin
Ðā dēor sindon bunden. Ðā dēor sindon bundene.
% mdparagraph:P160:end

% mdparagraph:P161:begin
The animals are bound. The animals are bound.
% mdparagraph:P161:end

% mdparagraph:P162:begin
Both are licensed verbal passives: bunden is uninflected and bundene has the masculine-like plural ending. The noun dēor remains neuter plural in both.
% mdparagraph:P162:end

% table: U07-DISPLAY-PREDICATE-PART
% mdtable:T37:begin
\par\begingroup
\fontsize{10.5}{13.2}\selectfont
\setlength{\LTleft}{0pt}\setlength{\LTright}{\fill}
\noindent
\begin{tabular}{@{}>{\raggedright\arraybackslash}p{38.41mm}>{\raggedright\arraybackslash}p{32.92mm}>{\raggedright\arraybackslash}p{47.05mm}@{}}
\toprule
\textbf{Verb} & \textbf{Uninflected} & \textbf{Plural masculine-like} \\
\midrule
bindan & bunden & bundene \\
gemanigfieldan & gemanigfield & gemanigfielde \\
\bottomrule
\end{tabular}
\par\endgroup

% mdtable:T37:end
% mdparagraph:P163:begin
The spellings can overlap across constructions. Attributive neuter plural bundenu has the accepted variants bunden and bundene given in §5. A spelling used in the passive is therefore not automatically forbidden before a noun. Identify the construction first; apply its agreement rule or permitted predicate inflection.
% mdparagraph:P163:end

% mdparagraph:P164:begin
Present participles also have uninflected predicate uses in the sources. The independent adjectival practice here uses the full attributive models. Do not infer an English-style progressive system merely from a form in -ende.
% mdparagraph:P164:end

% teaching: NEW-U07-GRAM-U07-PASSIVE-COMMAND-WISH
\reserve{7\baselineskip}
% mdparagraph:P165:begin
\subsection*{Passive commands and wishes}
% mdparagraph:P165:end

% mdparagraph:P166:begin
Use the imperative of the auxiliary for a direct command, and its present subjunctive for a wish. The participle remains nonfinite. In the first utterance below the speaker explicitly commands one male addressee; the second is a wish about the birds.
% mdparagraph:P166:end

% example: NEW-U07-EXAMPLE-10
% mdparagraph:P167:begin
Bēo geblētsod. Bēon ðā fuglas gemanigfielde.
% mdparagraph:P167:end

% mdparagraph:P168:begin
Be blessed. Let the birds be multiplied.
% mdparagraph:P168:end

% mdparagraph:P169:begin
Bēo is second-singular imperative in the stated command setting; bēon is third-plural subjunctive with ðā fuglas. Bēo alone could also be a subjunctive spelling, so the command analysis depends on the setting. The source-style plural gemanigfielde is a past participle, not a finite verb.
% mdparagraph:P169:end

% mdparagraph:P170:begin
This unit uses ordinary passives of the selected transitive verbs. It does not add agent phrases, a general perfect system, or a general progressive system. A participle of an intransitive verb may have a different function; the Creation closing formula is supplied locally in §10.
% mdparagraph:P170:end

\reserve{7\baselineskip}
% mdparagraph:P171:begin
\subsection*{Practice P. Participles and voice}
% mdparagraph:P171:end

% mdparagraph:P172:begin
Use the full adjective models where agreement is requested. For a verbal passive, distinguish required subject–auxiliary agreement from the permitted forms of its predicate participle.
% mdparagraph:P172:end

\reserve{5\baselineskip}
\par\begin{minipage}{\linewidth}\raggedright
% mdparagraph:P173:begin
\textbf{P1.}
% mdparagraph:P173:end

% item: NEW-U07-P01
% mdparagraph:P174:begin
Form the present participle of bindan, nęrian, lufian, and cweðan. Give the uninflected dictionary form and identify the stem retained before -ende.
% mdparagraph:P174:end

\end{minipage}\par\addvspace{4pt}

\par\begin{minipage}{\linewidth}\raggedright
% mdparagraph:P175:begin
\textbf{P2.}
% mdparagraph:P175:end

% item: NEW-U07-P02
% mdparagraph:P176:begin
Give the selected past participle of bindan, beran, hīeran, and gehālgian. Explain why the first two cannot be formed by adding a weak dental ending.
% mdparagraph:P176:end

\end{minipage}\par\addvspace{4pt}

\par\begin{minipage}{\linewidth}\raggedright
% mdparagraph:P177:begin
\textbf{P3.}
% mdparagraph:P177:end

% item: NEW-U07-P03
% mdparagraph:P178:begin
Supply the agreeing form of bindende: (a) a binding man, nominative; (b) the binding man, nominative; (c) with a binding woman, dative feminine singular only; (d) of the binding men, genitive plural only. For (c) and (d), give the participial form only; no new noun or preposition is required.
% mdparagraph:P178:end

\end{minipage}\par\addvspace{4pt}

\par\begin{minipage}{\linewidth}\raggedright
% mdparagraph:P179:begin
\textbf{P4.}
% mdparagraph:P179:end

% item: NEW-U07-P04
% mdparagraph:P180:begin
Supply the agreeing form of bunden: (a) a bound animal, accusative neuter singular; (b) the bound animal, accusative neuter singular; (c) of bound birds, genitive plural; (d) to the bound birds, dative plural. Identify strong or weak inflection in each.
% mdparagraph:P180:end

\end{minipage}\par\addvspace{4pt}

\par\begin{minipage}{\linewidth}\raggedright
% mdparagraph:P181:begin
\textbf{P5.}
% mdparagraph:P181:end

% item: NEW-U07-P05
% mdparagraph:P182:begin
Change Hē bint ðæt dēor into a present passive with wesan. Omit the agent. Identify the subject and the case of dēor before and after the change.
% mdparagraph:P182:end

\end{minipage}\par\addvspace{4pt}

\par\begin{minipage}{\linewidth}\raggedright
% mdparagraph:P183:begin
\textbf{P6.}
% mdparagraph:P183:end

% item: NEW-U07-P06
% mdparagraph:P184:begin
Change Hīe bundon ðā dēor into a past passive with weorðan, describing the binding event. Omit the agent.
% mdparagraph:P184:end

\end{minipage}\par\addvspace{4pt}

\par\begin{minipage}{\linewidth}\raggedright
% mdparagraph:P185:begin
\textbf{P7.}
% mdparagraph:P185:end

% item: NEW-U07-P07
% mdparagraph:P186:begin
Explain why both Ðā dēor sindon bunden and Ðā dēor sindon bundene can serve as verbal passives. Contrast these with an attributive participle directly modifying a neuter plural noun.
% mdparagraph:P186:end

\end{minipage}\par\addvspace{4pt}

\par\begin{minipage}{\linewidth}\raggedright
% mdparagraph:P187:begin
\textbf{P8.}
% mdparagraph:P187:end

% item: NEW-U07-P08
% mdparagraph:P188:begin
The first utterance is explicitly a command to one male addressee: Bēo geblētsod. The second is a wish about the birds: Bēon ðā fuglas gemanigfielde. Parse each auxiliary and participle, explain the addressee/\allowbreak{}subject difference, and translate both. Note any formal ambiguity that the stated context resolves.
% mdparagraph:P188:end

\end{minipage}\par\addvspace{4pt}

\reserve{7\baselineskip}
% mdparagraph:P189:begin
\section*{7. Quantifiers, places, and functions}
% mdparagraph:P189:end

% teaching: NEW-U07-VOC-eall-all
% teaching: NEW-U07-VOC-aelc-each
% teaching: NEW-U07-GRAM-U07-QUANTIFIER-AGREEMENT
\reserve{7\baselineskip}
% mdparagraph:P190:begin
\subsection*{Eall and ǣlc}
% mdparagraph:P190:end

% mdparagraph:P191:begin
\textbf{Eall}, “all; whole”, and \textbf{ǣlc}, “each; every”, agree with the noun they qualify. Learn their complete strong forms. Eall can refer to a whole singular thing or to all the members of a plural group. The ordinary “each/\allowbreak{}every” use of ǣlc in this unit takes a \textbf{singular noun}. Its full plural morphology is given, but no new plural meaning or partitive construction is required in independent translation.
% mdparagraph:P191:end

% teaching: NEW-U07-NOM-eall
\reserve{9\baselineskip}
% mdparagraph:P192:begin
\textbf{Eall: singular.}
% mdparagraph:P192:end

% table: NEW-U07-NOM-eall-1
% mdtable:T38:begin
% course-table-style: 2026-10
% course-panel:T38-P01:begin
\par\begingroup
\fontsize{10.5}{13.2}\selectfont
\setlength{\LTleft}{0pt}\setlength{\LTright}{\fill}
\noindent
\begin{tabular}{@{}>{\raggedright\arraybackslash}p{22.0mm}>{\raggedright\arraybackslash}p{31.188mm}>{\raggedright\arraybackslash}p{31.188mm}>{\raggedright\arraybackslash}p{31.188mm}@{}}
\toprule
\textbf{Case} & \textbf{Masculine} & \textbf{Neuter} & \textbf{Feminine} \\
\midrule
Nominative & eall & eall & eall \\
Genitive & ealles & ealles & ealre \\
Dative & eallum & eallum & ealre \\
Accusative & ealne & eall & ealle \\
\bottomrule
\end{tabular}
\par\endgroup
% course-panel:T38-P01:end
% mdtable:T38:end
\reserve{9\baselineskip}
% mdparagraph:P193:begin
\textbf{Eall: plural.}
% mdparagraph:P193:end

% table: NEW-U07-NOM-eall-2
% mdtable:T39:begin
% course-table-style: 2026-10
% course-panel:T39-P01:begin
\par\begingroup
\fontsize{10.5}{13.2}\selectfont
\setlength{\LTleft}{0pt}\setlength{\LTright}{\fill}
\noindent
\begin{tabular}{@{}>{\raggedright\arraybackslash}p{22.0mm}>{\raggedright\arraybackslash}p{31.188mm}>{\raggedright\arraybackslash}p{31.188mm}>{\raggedright\arraybackslash}p{31.188mm}@{}}
\toprule
\textbf{Case} & \textbf{Masculine} & \textbf{Neuter} & \textbf{Feminine} \\
\midrule
Nominative & ealle & eall & ealla \\
Genitive & ealra & ealra & ealra \\
Dative & eallum & eallum & eallum \\
Accusative & ealle & eall & ealla \\
\bottomrule
\end{tabular}
\par\endgroup
% course-panel:T39-P01:end
% mdtable:T39:end
% mdparagraph:P194:begin
Preserve the single l in ealne, ealre, and ealra, beside eall- in the other printed forms. These are agreement forms, not an invariant English “all”.
% mdparagraph:P194:end

% teaching: NEW-U07-NOM-aelc
\reserve{9\baselineskip}
% mdparagraph:P195:begin
\textbf{Ǣlc: singular.}
% mdparagraph:P195:end

% table: NEW-U07-NOM-aelc-1
% mdtable:T40:begin
% course-table-style: 2026-10
% course-panel:T40-P01:begin
\par\begingroup
\fontsize{10.5}{13.2}\selectfont
\setlength{\LTleft}{0pt}\setlength{\LTright}{\fill}
\noindent
\begin{tabular}{@{}>{\raggedright\arraybackslash}p{22.0mm}>{\raggedright\arraybackslash}p{31.188mm}>{\raggedright\arraybackslash}p{31.188mm}>{\raggedright\arraybackslash}p{31.188mm}@{}}
\toprule
\textbf{Case} & \textbf{Masculine} & \textbf{Neuter} & \textbf{Feminine} \\
\midrule
Nominative & ǣlc & ǣlc & ǣlc \\
Genitive & ǣlces & ǣlces & ǣlcre \\
Dative & ǣlcum & ǣlcum & ǣlcre \\
Accusative & ǣlcne & ǣlc & ǣlce \\
\bottomrule
\end{tabular}
\par\endgroup
% course-panel:T40-P01:end
% mdtable:T40:end
\reserve{9\baselineskip}
% mdparagraph:P196:begin
\textbf{Ǣlc: plural morphology.}
% mdparagraph:P196:end

% table: NEW-U07-NOM-aelc-2
% mdtable:T41:begin
% course-table-style: 2026-10
% course-panel:T41-P01:begin
\par\begingroup
\fontsize{10.5}{13.2}\selectfont
\setlength{\LTleft}{0pt}\setlength{\LTright}{\fill}
\noindent
\begin{tabular}{@{}>{\raggedright\arraybackslash}p{22.0mm}>{\raggedright\arraybackslash}p{31.188mm}>{\raggedright\arraybackslash}p{31.188mm}>{\raggedright\arraybackslash}p{31.188mm}@{}}
\toprule
\textbf{Case} & \textbf{Masculine} & \textbf{Neuter} & \textbf{Feminine} \\
\midrule
Nominative & ǣlce & ǣlc & ǣlca \\
Genitive & ǣlcra & ǣlcra & ǣlcra \\
Dative & ǣlcum & ǣlcum & ǣlcum \\
Accusative & ǣlce & ǣlc & ǣlca \\
\bottomrule
\end{tabular}
\par\endgroup
% course-panel:T41-P01:end
% mdtable:T41:end
% mdparagraph:P197:begin
The vowel is long ǣ throughout. Use the singular case required by the noun’s role when translating “each day” or “every bird”. Displaying the plural endings does not change that ordinary construction. These are strong paradigms; no new weak quantifier system or instrumental use is introduced.
% mdparagraph:P197:end

\reserve{7\baselineskip}
% mdparagraph:P198:begin
\subsection*{Geweald: dominion}
% mdparagraph:P198:end

% teaching: NEW-U07-NOM-geweald
% teaching: NEW-U07-VOC-geweald-dominion
% mdparagraph:P199:begin
\textbf{Geweald} is a strong neuter noun, “control; rule; dominion”. Use the singular for the ordinary meaning in this unit. Its complete declension includes the plural for recognition and morphological reference; translating independent plural “dominions” is not required.
% mdparagraph:P199:end

% table: NEW-U07-NOM-geweald
% mdtable:T42:begin
% course-table-style: 2026-10
% course-panel:T42-P01:begin
\par\begingroup
\fontsize{10.5}{13.2}\selectfont
\setlength{\LTleft}{0pt}\setlength{\LTright}{\fill}
\noindent
\begin{tabular}{@{}>{\raggedright\arraybackslash}p{39.459mm}>{\raggedright\arraybackslash}p{39.459mm}>{\raggedright\arraybackslash}p{39.459mm}@{}}
\toprule
\textbf{Case} & \textbf{Singular} & \textbf{Plural (morphology reference)} \\
\midrule
Nominative & geweald & geweald \\
Genitive & gewealdes & gewealda \\
Dative & gewealde & gewealdum \\
Accusative & geweald & geweald \\
\bottomrule
\end{tabular}
\par\endgroup
% course-panel:T42-P01:end
% mdtable:T42:end
% mdparagraph:P200:begin
The long neuter stem has the same nominative and accusative spelling in singular and plural. The genitive singular is -es and dative singular -e; genitive plural -a and dative plural -um distinguish the oblique plural forms.
% mdparagraph:P200:end

\reserve{7\baselineskip}
% mdparagraph:P201:begin
\subsection*{Bearing, location, and assigned function}
% mdparagraph:P201:end

% teaching: NEW-U07-VOC-beran-produce
% teaching: NEW-U07-GRAM-U07-BERAN-ACCUSATIVE-SEED-FRUIT
% mdparagraph:P202:begin
\textbf{Beran} can mean “bear” in the sense \textbf{produce seed or fruit}, with the seed or fruit as accusative object. This use is additional to the earlier literal “carry”. Wyrcean “yield fruit” likewise takes an accusative fruit object. Keep the short vowel of wæstm, “fruit”.
% mdparagraph:P202:end

% teaching: NEW-U07-VOC-ofer
% teaching: NEW-U07-VOC-ofer-authority
% teaching: NEW-U07-VOC-ofer-1
% teaching: NEW-U07-VOC-bufan
% teaching: NEW-U07-VOC-bufan-above
% teaching: NEW-U07-VOC-under
% teaching: NEW-U07-VOC-under-below
% teaching: NEW-U07-VOC-tomiddes
% teaching: NEW-U07-VOC-tomiddes-midst
% teaching: NEW-U07-VOC-to
% teaching: NEW-U07-VOC-to-function-purpose
% teaching: NEW-U07-GRAM-U07-OFER-ACCUSATIVE
% teaching: NEW-U07-GRAM-U07-SPATIAL-OFER
% teaching: NEW-U07-GRAM-U07-BUFAN-DATIVE
% teaching: NEW-U07-GRAM-U07-UNDER-DATIVE
% teaching: NEW-U07-GRAM-U07-TOMIDDES-DATIVE
% teaching: NEW-U07-GRAM-U07-TO-DATIVE
% teaching: NEW-U07-GRAM-U07-DATIVE-FUNCTION-BENEFICIARY
% mdparagraph:P203:begin
The following prepositions have the stated meanings and cases here. A familiar spelling does not make all its uses interchangeable.
% mdparagraph:P203:end

% table: U07-DISPLAY-FUNCTION-WORDS
% mdtable:T43:begin
\par\begingroup
\fontsize{10.5}{13.2}\selectfont
\setlength{\LTleft}{0pt}\setlength{\LTright}{\fill}
\noindent
\begin{tabular}{@{}>{\raggedright\arraybackslash}p{21.58mm}>{\raggedright\arraybackslash}p{33.02mm}>{\raggedright\arraybackslash}p{63.77mm}@{}}
\toprule
\textbf{Word} & \textbf{Meaning here} & \textbf{Construction} \\
\midrule
ofer & over, in authority over & accusative in the authority frame \\
ofer & spatially over, across & accusative or dative, according to the construction \\
bufan & above & dative landmark \\
under & under, below & dative landmark in the location use taught here \\
tōmiddes & in the midst of & dative landmark \\
tō & for, serving as & dative assigned function, such as food or a sign \\
\bottomrule
\end{tabular}
\par\endgroup

% mdtable:T43:end
% mdparagraph:P204:begin
With spatial ofer, the accusative commonly marks direction over or across a landmark, while the dative commonly marks location. This is a tendency, not an exceptionless rule for every use. The source’s ofer eorðan cannot by its spelling alone identify accusative rather than dative: eorðan serves both cases. The authority use is separately taught with the accusative.
% mdparagraph:P204:end

% mdparagraph:P205:begin
Tō with a dative noun can state what something is \textbf{for} or \textbf{serves as}. A separate dative can name the beneficiary: in ēow tō męte, “to you for food”, ēow is the beneficiary and męte is the assigned function. Neither is a possessive adjective. This tō is different from the earlier tō with an inflected purpose infinitive.
% mdparagraph:P205:end

\reserve{7\baselineskip}
% mdparagraph:P206:begin
\subsection*{Naming}
% mdparagraph:P206:end

% teaching: NEW-U07-GRAM-U07-NAMING-TWO-ACCUSATIVES
% mdparagraph:P207:begin
Hātan “call, name” can take \textbf{two accusatives}: the thing named and the name given to it. In the first sentence below the firmament is the thing named and Heaven the assigned name. The second sentence shows tō with a dative function in a wish.
% mdparagraph:P207:end

% example: NEW-U07-EXAMPLE-11
% mdparagraph:P208:begin
Ic hāte ðā fæstnesse Heofon. Bēo ðæt lēoht tō tācne.
% mdparagraph:P208:end

% mdparagraph:P209:begin
I call the firmament Heaven. Let the light be for a sign.
% mdparagraph:P209:end

% mdparagraph:P210:begin
Ðā fæstnesse and Heofon are both accusative. In the wish, ðæt lēoht is nominative subject and tācne is dative after tō. The “be called” passive uses the learned auxiliary with hāten; a separate hātte system is not needed for this unit.
% mdparagraph:P210:end

\reserve{7\baselineskip}
% mdparagraph:P211:begin
\subsection*{Adverbs and a presentational word}
% mdparagraph:P211:end

% teaching: NEW-U07-VOC-nu
% teaching: NEW-U07-VOC-nu-1
% teaching: NEW-U07-VOC-swa
% teaching: NEW-U07-VOC-swa-1
% teaching: NEW-U07-VOC-thus
% teaching: NEW-U07-VOC-thus-thus-source
% teaching: NEW-U07-VOC-eac
% teaching: NEW-U07-VOC-eac-also
% teaching: NEW-U07-VOC-swilce
% teaching: NEW-U07-VOC-swilce-likewise
% teaching: NEW-U07-VOC-efne
% teaching: NEW-U07-VOC-efne-behold
% teaching: NEW-U07-GRAM-U07-ADVERBS-INTERJECTION
% mdtable:T44:begin
\par\begingroup
\fontsize{10.5}{13.2}\selectfont
\setlength{\LTleft}{0pt}\setlength{\LTright}{\fill}
\noindent
\begin{tabular}{@{}>{\raggedright\arraybackslash}p{25.49mm}>{\raggedright\arraybackslash}p{95.7mm}@{}}
\toprule
\textbf{Word} & \textbf{Meaning and use} \\
\midrule
nū & now; can open a statement with is before its following nominal subject \\
swā & so, in that way \\
ðus & thus, in this way \\
ēac & also \\
swilce & likewise; the adverb, not a newly taught conjunction \\
efne & behold; introduces the following statement \\
gif & if; introduces the condition described in §4 \\
\bottomrule
\end{tabular}
\par\endgroup

% mdtable:T44:end
% mdparagraph:P212:begin
All these forms are invariant. The source also uses ēac swilce, “also”. The already learned swīðe means “very” before an adjective. Keep the adjective’s own agreement: the intensifier does not alter it.
% mdparagraph:P212:end

% mdparagraph:P213:begin
\textbf{Predicate nouns.} The finite verb agrees with its subject. A predicate noun may have its own number: in “Birds are the food of the animal”, the birds are plural but “food” is singular męte. The rule requiring an adjective to agree in number does not turn a singular predicate noun into a plural.
% mdparagraph:P213:end

\reserve{7\baselineskip}
% mdparagraph:P214:begin
\section*{8. Translation practice}
% mdparagraph:P214:end

\reserve{7\baselineskip}
% mdparagraph:P215:begin
\subsection*{Practice C. Controlled composition}
% mdparagraph:P215:end

% mdparagraph:P216:begin
Use the stated command, wish, or conditional target. Retain the case required by the verb or preposition. The last item contrasts singular “every” with plural “all”.
% mdparagraph:P216:end

\reserve{5\baselineskip}
\par\begin{minipage}{\linewidth}\raggedright
% mdparagraph:P217:begin
\textbf{C1.}
% mdparagraph:P217:end

% item: NEW-U07-C01
% mdparagraph:P218:begin
Translate: May the birds live.
% mdparagraph:P218:end

\end{minipage}\par\addvspace{4pt}

\par\begin{minipage}{\linewidth}\raggedright
% mdparagraph:P219:begin
\textbf{C2.}
% mdparagraph:P219:end

% item: NEW-U07-C02
% mdparagraph:P220:begin
Translate: Let us help the companions.
% mdparagraph:P220:end

\end{minipage}\par\addvspace{4pt}

\par\begin{minipage}{\linewidth}\raggedright
% mdparagraph:P221:begin
\textbf{C3.}
% mdparagraph:P221:end

% item: NEW-U07-C03
% mdparagraph:P222:begin
Translate, presenting the condition as tentative: If you (sing.) have food, give it to the birds.
% mdparagraph:P222:end

\end{minipage}\par\addvspace{4pt}

\par\begin{minipage}{\linewidth}\raggedright
% mdparagraph:P223:begin
\textbf{C4.}
% mdparagraph:P223:end

% item: NEW-U07-C04
% mdparagraph:P224:begin
Translate as a command to one male addressee: Be blessed.
% mdparagraph:P224:end

\end{minipage}\par\addvspace{4pt}

\par\begin{minipage}{\linewidth}\raggedright
% mdparagraph:P225:begin
\textbf{C5.}
% mdparagraph:P225:end

% item: NEW-U07-C05
% mdparagraph:P226:begin
Translate as a wish: Let the work be accomplished.
% mdparagraph:P226:end

\end{minipage}\par\addvspace{4pt}

\par\begin{minipage}{\linewidth}\raggedright
% mdparagraph:P227:begin
\textbf{C6.}
% mdparagraph:P227:end

% item: NEW-U07-C06
% mdparagraph:P228:begin
Translate: Every bird is above the earth; all trees bear fruit.
% mdparagraph:P228:end

\end{minipage}\par\addvspace{4pt}

\reserve{7\baselineskip}
% mdparagraph:P229:begin
\subsection*{Cook’s exercises}
% mdparagraph:P229:end

% mdparagraph:P230:begin
Translate into Old English. The original wording and numbering are retained. For English “let” or “may” expressing a wish, use the subjunctive; direct commands use the imperative. Use the taught auxiliary–participle construction for passives and the naming construction for “call”. Keep direct speech as direct speech in I.7.17.
% mdparagraph:P230:end

% mdparagraph:P231:begin
Where English “you” has no number specified, either singular or plural is possible; make the pronouns and related forms consistent. The parenthetical number directions in the original prompts do specify a target. An otherwise valid translation can differ from a model’s chosen wording or number.
% mdparagraph:P231:end

\reserve{16\baselineskip}
% mdparagraph:P232:begin
\subsection*{Cook I.4: the remaining item}
% mdparagraph:P232:end

\reserve{5\baselineskip}
\par\begin{minipage}{\linewidth}\raggedright
% mdparagraph:P233:begin
\textbf{Cook I.4.18.}
% mdparagraph:P233:end

% item: EX-I-04-18
% mdparagraph:P234:begin
Increase (sing.) and be multiplied.
% mdparagraph:P234:end

\end{minipage}\par\addvspace{4pt}

\reserve{7\baselineskip}
% mdparagraph:P235:begin
\subsection*{Cook I.6}
% mdparagraph:P235:end

\reserve{5\baselineskip}
\par\begin{minipage}{\linewidth}\raggedright
% mdparagraph:P236:begin
\textbf{Cook I.6.1.}
% mdparagraph:P236:end

% item: EX-I-06-01
% mdparagraph:P237:begin
The animal lives.
% mdparagraph:P237:end

\end{minipage}\par\addvspace{4pt}

\par\begin{minipage}{\linewidth}\raggedright
% mdparagraph:P238:begin
\textbf{Cook I.6.2.}
% mdparagraph:P238:end

% item: EX-I-06-02
% mdparagraph:P239:begin
The bird flies.
% mdparagraph:P239:end

\end{minipage}\par\addvspace{4pt}

\par\begin{minipage}{\linewidth}\raggedright
% mdparagraph:P240:begin
\textbf{Cook I.6.3.}
% mdparagraph:P240:end

% item: EX-I-06-03
% mdparagraph:P241:begin
The tree grows.
% mdparagraph:P241:end

\end{minipage}\par\addvspace{4pt}

\par\begin{minipage}{\linewidth}\raggedright
% mdparagraph:P242:begin
\textbf{Cook I.6.4.}
% mdparagraph:P242:end

% item: EX-I-06-04
% mdparagraph:P243:begin
The man rules over the cattle of the earth.
% mdparagraph:P243:end

\end{minipage}\par\addvspace{4pt}

\par\begin{minipage}{\linewidth}\raggedright
% mdparagraph:P244:begin
\textbf{Cook I.6.5.}
% mdparagraph:P244:end

% item: EX-I-06-05
% mdparagraph:P245:begin
Let him be over the birds and (over) the beasts.
% mdparagraph:P245:end

\end{minipage}\par\addvspace{4pt}

\par\begin{minipage}{\linewidth}\raggedright
% mdparagraph:P246:begin
\textbf{Cook I.6.6.}
% mdparagraph:P246:end

% item: EX-I-06-06
% mdparagraph:P247:begin
The stars are above the sea.
% mdparagraph:P247:end

\end{minipage}\par\addvspace{4pt}

\par\begin{minipage}{\linewidth}\raggedright
% mdparagraph:P248:begin
\textbf{Cook I.6.7.}
% mdparagraph:P248:end

% item: EX-I-06-07
% mdparagraph:P249:begin
It is for your food [to you for food].
% mdparagraph:P249:end

\end{minipage}\par\addvspace{4pt}

\par\begin{minipage}{\linewidth}\raggedright
% mdparagraph:P250:begin
\textbf{Cook I.6.8.}
% mdparagraph:P250:end

% item: EX-I-06-08
% mdparagraph:P251:begin
The herb bears seed.
% mdparagraph:P251:end

\end{minipage}\par\addvspace{4pt}

\par\begin{minipage}{\linewidth}\raggedright
% mdparagraph:P252:begin
\textbf{Cook I.6.9.}
% mdparagraph:P252:end

% item: EX-I-06-09
% mdparagraph:P253:begin
The tree yields fruit.
% mdparagraph:P253:end

\end{minipage}\par\addvspace{4pt}

\par\begin{minipage}{\linewidth}\raggedright
% mdparagraph:P254:begin
\textbf{Cook I.6.10.}
% mdparagraph:P254:end

% item: EX-I-06-10
% mdparagraph:P255:begin
The air is very good.
% mdparagraph:P255:end

\end{minipage}\par\addvspace{4pt}

\par\begin{minipage}{\linewidth}\raggedright
% mdparagraph:P256:begin
\textbf{Cook I.6.11.}
% mdparagraph:P256:end

% item: EX-I-06-11
% mdparagraph:P257:begin
The things are very good.
% mdparagraph:P257:end

\end{minipage}\par\addvspace{4pt}

\par\begin{minipage}{\linewidth}\raggedright
% mdparagraph:P258:begin
\textbf{Cook I.6.12.}
% mdparagraph:P258:end

% item: EX-I-06-12
% mdparagraph:P259:begin
The birds were very good.
% mdparagraph:P259:end

\end{minipage}\par\addvspace{4pt}

\par\begin{minipage}{\linewidth}\raggedright
% mdparagraph:P260:begin
\textbf{Cook I.6.13.}
% mdparagraph:P260:end

% item: EX-I-06-13
% mdparagraph:P261:begin
The light of the stars was very good.
% mdparagraph:P261:end

\end{minipage}\par\addvspace{4pt}

\par\begin{minipage}{\linewidth}\raggedright
% mdparagraph:P262:begin
\textbf{Cook I.6.14.}
% mdparagraph:P262:end

% item: EX-I-06-14
% mdparagraph:P263:begin
He gave them all trees.
% mdparagraph:P263:end

\end{minipage}\par\addvspace{4pt}

\par\begin{minipage}{\linewidth}\raggedright
% mdparagraph:P264:begin
\textbf{Cook I.6.15.}
% mdparagraph:P264:end

% item: EX-I-06-15
% mdparagraph:P265:begin
I give you all things upon earth.
% mdparagraph:P265:end

\end{minipage}\par\addvspace{4pt}

\par\begin{minipage}{\linewidth}\raggedright
% mdparagraph:P266:begin
\textbf{Cook I.6.16.}
% mdparagraph:P266:end

% item: EX-I-06-16
% mdparagraph:P267:begin
I call the firmament Heaven.
% mdparagraph:P267:end

\end{minipage}\par\addvspace{4pt}

\reserve{7\baselineskip}
% mdparagraph:P268:begin
\subsection*{Cook I.7}
% mdparagraph:P268:end

\reserve{5\baselineskip}
\par\begin{minipage}{\linewidth}\raggedright
% mdparagraph:P269:begin
\textbf{Cook I.7.1.}
% mdparagraph:P269:end

% item: EX-I-07-01
% mdparagraph:P270:begin
He has dominion.
% mdparagraph:P270:end

\end{minipage}\par\addvspace{4pt}

\par\begin{minipage}{\linewidth}\raggedright
% mdparagraph:P271:begin
\textbf{Cook I.7.2.}
% mdparagraph:P271:end

% item: EX-I-07-02
% mdparagraph:P272:begin
He blesses the earth.
% mdparagraph:P272:end

\end{minipage}\par\addvspace{4pt}

\par\begin{minipage}{\linewidth}\raggedright
% mdparagraph:P273:begin
\textbf{Cook I.7.3.}
% mdparagraph:P273:end

% item: EX-I-07-03
% mdparagraph:P274:begin
He says thus.
% mdparagraph:P274:end

\end{minipage}\par\addvspace{4pt}

\par\begin{minipage}{\linewidth}\raggedright
% mdparagraph:P275:begin
\textbf{Cook I.7.4.}
% mdparagraph:P275:end

% item: EX-I-07-04
% mdparagraph:P276:begin
Now is the season.
% mdparagraph:P276:end

\end{minipage}\par\addvspace{4pt}

\par\begin{minipage}{\linewidth}\raggedright
% mdparagraph:P277:begin
\textbf{Cook I.7.5.}
% mdparagraph:P277:end

% item: EX-I-07-05
% mdparagraph:P278:begin
It was so called.
% mdparagraph:P278:end

\end{minipage}\par\addvspace{4pt}

\par\begin{minipage}{\linewidth}\raggedright
% mdparagraph:P279:begin
\textbf{Cook I.7.6.}
% mdparagraph:P279:end

% item: EX-I-07-06
% mdparagraph:P280:begin
Behold, God rules over the earth.
% mdparagraph:P280:end

\end{minipage}\par\addvspace{4pt}

\par\begin{minipage}{\linewidth}\raggedright
% mdparagraph:P281:begin
\textbf{Cook I.7.7.}
% mdparagraph:P281:end

% item: EX-I-07-07
% mdparagraph:P282:begin
The tree is in the midst of the grass.
% mdparagraph:P282:end

\end{minipage}\par\addvspace{4pt}

\par\begin{minipage}{\linewidth}\raggedright
% mdparagraph:P283:begin
\textbf{Cook I.7.8.}
% mdparagraph:P283:end

% item: EX-I-07-08
% mdparagraph:P284:begin
The earth is in the midst of the air.
% mdparagraph:P284:end

\end{minipage}\par\addvspace{4pt}

\par\begin{minipage}{\linewidth}\raggedright
% mdparagraph:P285:begin
\textbf{Cook I.7.9.}
% mdparagraph:P285:end

% item: EX-I-07-09
% mdparagraph:P286:begin
The waters are under the firmament.
% mdparagraph:P286:end

\end{minipage}\par\addvspace{4pt}

\par\begin{minipage}{\linewidth}\raggedright
% mdparagraph:P287:begin
\textbf{Cook I.7.10.}
% mdparagraph:P287:end

% item: EX-I-07-10
% mdparagraph:P288:begin
Birds are the food of the animal.
% mdparagraph:P288:end

\end{minipage}\par\addvspace{4pt}

\par\begin{minipage}{\linewidth}\raggedright
% mdparagraph:P289:begin
\textbf{Cook I.7.11.}
% mdparagraph:P289:end

% item: EX-I-07-11
% mdparagraph:P290:begin
Divide (plur.) the waters from the waters.
% mdparagraph:P290:end

\end{minipage}\par\addvspace{4pt}

\par\begin{minipage}{\linewidth}\raggedright
% mdparagraph:P291:begin
\textbf{Cook I.7.12.}
% mdparagraph:P291:end

% item: EX-I-07-12
% mdparagraph:P292:begin
God's dominion is over all creatures.
% mdparagraph:P292:end

\end{minipage}\par\addvspace{4pt}

\par\begin{minipage}{\linewidth}\raggedright
% mdparagraph:P293:begin
\textbf{Cook I.7.13.}
% mdparagraph:P293:end

% item: EX-I-07-13
% mdparagraph:P294:begin
God hallows each day.
% mdparagraph:P294:end

\end{minipage}\par\addvspace{4pt}

\par\begin{minipage}{\linewidth}\raggedright
% mdparagraph:P295:begin
\textbf{Cook I.7.14.}
% mdparagraph:P295:end

% item: EX-I-07-14
% mdparagraph:P296:begin
Let the birds increase and be multiplied.
% mdparagraph:P296:end

\end{minipage}\par\addvspace{4pt}

\par\begin{minipage}{\linewidth}\raggedright
% mdparagraph:P297:begin
\textbf{Cook I.7.15.}
% mdparagraph:P297:end

% item: EX-I-07-15
% mdparagraph:P298:begin
Let the men swim in the sea.
% mdparagraph:P298:end

\end{minipage}\par\addvspace{4pt}

\par\begin{minipage}{\linewidth}\raggedright
% mdparagraph:P299:begin
\textbf{Cook I.7.16.}
% mdparagraph:P299:end

% item: EX-I-07-16
% mdparagraph:P300:begin
Let the tree bear fruit.
% mdparagraph:P300:end

\end{minipage}\par\addvspace{4pt}

\par\begin{minipage}{\linewidth}\raggedright
% mdparagraph:P301:begin
\textbf{Cook I.7.17.}
% mdparagraph:P301:end

% item: EX-I-07-17
% mdparagraph:P302:begin
He says likewise, Let the light be for a sign.
% mdparagraph:P302:end

\end{minipage}\par\addvspace{4pt}

\reserve{7\baselineskip}
% mdparagraph:P303:begin
\section*{9. Binding the animals}
% mdparagraph:P303:end

% teaching: NEW-U07-GRAM-U07-QUOTATION-WISH
% mdparagraph:P304:begin
\textbf{New composition.} The two speeches distinguish a tentative condition, an assertion of ability, willingness to help, and a joint exhortation. Pronouns inside each quotation belong to that speaker’s setting.
% mdparagraph:P304:end

% composition: NEW-U07-TEXT-A
% mdparagraph:P305:begin
\par\begin{minipage}{\linewidth}\raggedright
\readingtext{Twēgen męnn wǣron on ðǣm strande. Hīe hæfdon fīf dēor. Se ǣresta mann cwæð: “Gif ðū mē helpe, wē magon ðā dēor bindan.” Se ōðer mann cwæð: “Ic wille ðē helpan. Binden wē ðā dēor.” Hīe bundon ðā dēor, and se ǣresta mann sealde his gefēran męte. Ðā dēor wǣron bunden, and ðā męnn hæfdon męte.}
\end{minipage}\par
% mdparagraph:P305:end

\reserve{7\baselineskip}
% mdparagraph:P306:begin
\subsection*{Practice D. The connected account}
% mdparagraph:P306:end

% mdparagraph:P307:begin
Read the account as a whole before identifying individual forms. Use the sequence of actions and speakers to establish the participants referred to.
% mdparagraph:P307:end

\reserve{5\baselineskip}
\par\begin{minipage}{\linewidth}\raggedright
% mdparagraph:P308:begin
\textbf{D1.}
% mdparagraph:P308:end

% item: NEW-U07-D01
% mdparagraph:P309:begin
Translate the complete new account “Binding the animals.” Keep the speakers and the final result clear.
% mdparagraph:P309:end

\end{minipage}\par\addvspace{4pt}

\par\begin{minipage}{\linewidth}\raggedright
% mdparagraph:P310:begin
\textbf{D2.}
% mdparagraph:P310:end

% item: NEW-U07-D02
% mdparagraph:P311:begin
In the first speech, identify the condition and the claimed ability. In the second, distinguish the speaker’s willingness from the joint exhortation.
% mdparagraph:P311:end

\end{minipage}\par\addvspace{4pt}

\par\begin{minipage}{\linewidth}\raggedright
% mdparagraph:P312:begin
\textbf{D3.}
% mdparagraph:P312:end

% item: NEW-U07-D03
% mdparagraph:P313:begin
Identify the cases of mē, ðē, and his gefēran in the account. Explain the two dative relations, and identify the possessor intended by his from the stated speaker sequence.
% mdparagraph:P313:end

\end{minipage}\par\addvspace{4pt}

\par\begin{minipage}{\linewidth}\raggedright
% mdparagraph:P314:begin
\textbf{D4.}
% mdparagraph:P314:end

% item: NEW-U07-D04
% mdparagraph:P315:begin
Identify the final passive. What earlier event establishes its result, and which participants performed that event? Explain how the narrative supplies this information even though the passive clause has no agent phrase.
% mdparagraph:P315:end

\end{minipage}\par\addvspace{4pt}

\reserve{7\baselineskip}
% mdparagraph:P316:begin
\section*{10. The fifth day of Creation}
% mdparagraph:P316:end

% mdparagraph:P317:begin
This complete paragraph is from Cook’s \emph{First Book}, Reader I, printed p.126 /\allowbreak{} one-based PDF p.142, editorial prose lines 5–14. The printed page has marginal line anchors; the consecutive prose-line count is editorial. The paragraph begins with God’s direction to the waters, describes the creatures and blessing, and closes the fifth day.
% mdparagraph:P317:end

% source: FB-I-CREATION-FIFTH-DAY-01
% mdparagraph:P318:begin
\par\begin{minipage}{\linewidth}\raggedright
\readingtext{God cwæð ēac swilce, “Tēon nū ðā wæteru forð swimmendu cynn cucu on līfe, and flēogendu cynn ofer eorðan under ðǣre heofonan fæstnesse.” And God gescēop ðā ðā miclan hwalas, and eall libbendu fisccynn and styriendlicu, ðe ðā wæteru tugon forð on hiera hīwum, and eall flēogendu cynn æfter hiera cynne; God geseah ðā ðæt hit gōd wæs. And blētsode hīe, ðus cweðende, “Weaxað, and bēoð gemanigfielde, and gefyllað ðǣre sǣ wæteru, and ðā fuglas bēon gemanigfielde ofer eorðan.” And ðā wæs geworden ǣfen and męrgen se fīfta dæg.}
\end{minipage}\par
% mdparagraph:P318:end

\reserve{7\baselineskip}
% mdparagraph:P319:begin
\subsection*{Forms for this reading}
% mdparagraph:P319:end

% teaching: NEW-U07-SOURCE-FORMS
% mdparagraph:P320:begin
Recognize the following three \textbf{exact present-participle forms}. Each is strong neuter accusative plural with its noun. Their -u forms are retained from the source. The assignment is recognition of these forms; no independent recall of other participle cells of these three verbs is required.
% mdparagraph:P320:end

% table: U07-DISPLAY-SOURCE-SUPPORT
% mdtable:T45:begin
\par\begingroup
\fontsize{10.5}{13.2}\selectfont
\setlength{\LTleft}{0pt}\setlength{\LTright}{\fill}
\noindent
\begin{tabular}{@{}>{\raggedright\arraybackslash}p{21.01mm}>{\raggedright\arraybackslash}p{19.01mm}>{\raggedright\arraybackslash}p{19.01mm}>{\raggedright\arraybackslash}p{56.54mm}@{}}
\toprule
\textbf{Source form} & \textbf{Verb} & \textbf{Meaning} & \textbf{Agreement in this passage} \\
\midrule
swimmendu & swimman & swimming & neuter accusative plural with cynn \\
flēogendu & flēogan & flying & neuter accusative plural with cynn, at both occurrences \\
libbendu & libban & living & neuter accusative plural with fisccynn \\
\bottomrule
\end{tabular}
\par\endgroup

% mdtable:T45:end
% mdparagraph:P321:begin
The following exact forms are \textbf{supplied help}, available throughout the reading tasks. They do not add independent vocabulary or construction practice. Their phrases are explained in the ten groups below.
% mdparagraph:P321:end

% mdtable:T46:begin
\par\begingroup
\fontsize{10.5}{13.2}\selectfont
\setlength{\LTleft}{0pt}\setlength{\LTright}{\fill}
\noindent
\begin{tabular}{@{}>{\raggedright\arraybackslash}p{27.91mm}>{\raggedright\arraybackslash}p{93.28mm}@{}}
\toprule
\textbf{Source form} & \textbf{Local analysis and meaning} \\
\midrule
forð & adverb “forth”; part of the supplied “bring forth” expression \\
cucu & irregular adjective, neuter accusative plural, “alive” \\
heofonan & genitive singular of weak feminine heofone, “heaven” \\
ðā & temporal adverb “then” where identified below; distinct from the article \\
fisccynn & neuter accusative plural, “fish-kinds” \\
styriendlicu & strong neuter accusative plural, “moving”; this source’s -u variant \\
ðe & relative pronoun, object “which” in the supplied clause \\
hīwum & dative plural of neuter hīw, “forms; kinds” \\
æfter & “according to”, with dative cynne in the supplied phrase \\
ðæt & conjunction “that” introducing the supplied content clause \\
blētsode & third-singular preterite indicative of plain blētsian, “blessed” \\
sǣ & feminine genitive singular, “of the sea”, in ðǣre sǣ wæteru \\
geworden & past participle in the supplied “come to be” closing formula \\
męrgen & nominative singular source spelling of “morning” \\
\bottomrule
\end{tabular}
\par\endgroup

% mdtable:T46:end
% mdparagraph:P322:begin
Gemanigfielde is the learned plural predicate past participle, used twice in the blessing. Its source analysis is nominative plural. Cweðende has the learned present-participle form; the accompanying-action construction in which it occurs still receives local help.
% mdparagraph:P322:end

\reserve{7\baselineskip}
% mdparagraph:P323:begin
\subsection*{Local help}
% mdparagraph:P323:end

% mdparagraph:P324:begin
Use all ten groups during the reading exercises. They supply the unfamiliar lexical meanings and larger constructions; the learned finite verbs, participles, and agreement can still be identified inside them.
% mdparagraph:P324:end

\reserve{5\baselineskip}
% mdparagraph:P325:begin
\textbf{1. Tēon … forð /\allowbreak{} tugon forð.}
% mdparagraph:P325:end

% support: U07-SRC-HELP01
% teaching: NEW-U07-VOC-teon-bring-forth-source
% teaching: NEW-U07-VOC-forth-forth-source
% teaching: NEW-U07-GRAM-U07-SOURCE-BRING-FORTH
% mdparagraph:P326:begin
\textbf{Bring forth /\allowbreak{} brought forth.} Tēon is present third-plural subjunctive with ðā wæteru as its subject. Tugon is past third-plural indicative. The split “bring forth” expression and this sense are supplied: the earlier “draw, pull” does not by itself give this meaning.
% mdparagraph:P326:end

\reserve{5\baselineskip}
% mdparagraph:P327:begin
\textbf{2. swimmendu cynn cucu on līfe; flēogendu cynn.}
% mdparagraph:P327:end

% support: U07-SRC-HELP02
% teaching: NEW-U07-VOC-cucu-alive-source
% teaching: NEW-U07-VOC-on-alive-source
% teaching: NEW-U07-GRAM-U07-SOURCE-ALIVE
% mdparagraph:P328:begin
\textbf{Swimming kinds, alive/\allowbreak{}living; flying kinds.} Cynn is neuter accusative plural, the things to be brought forth. The exact participles swimmendu and flēogendu agree with it. Cucu is the supplied irregular neuter accusative plural adjective “alive”; keep cucu on līfe together as the supplied alive expression. Do not infer that every -ende participle must mechanically lose e before -u.
% mdparagraph:P328:end

\reserve{5\baselineskip}
% mdparagraph:P329:begin
\textbf{3. ofer eorðan; under ðǣre heofonan fæstnesse.}
% mdparagraph:P329:end

% support: U07-SRC-HELP03
% teaching: NEW-U07-VOC-heofone-1
% mdparagraph:P330:begin
\textbf{Over the earth; under the firmament of heaven.} Heofonan is the genitive singular of the source’s weak feminine heofone. Fæstnesse is dative after under. The spelling ðǣre can be feminine genitive or dative singular: it may agree with genitive heofonan or with dative fæstnesse. The wording does not require a unique choice between those article attachments.
% mdparagraph:P330:end

% mdparagraph:P331:begin
Eorðan after spatial ofer likewise has the same spelling in accusative and dative. The sense is clear, but the letters alone do not prove a unique case at either occurrence.
% mdparagraph:P331:end

\reserve{5\baselineskip}
% mdparagraph:P332:begin
\textbf{4. ðā ðā miclan hwalas; eall libbendu fisccynn and styriendlicu.}
% mdparagraph:P332:end

% support: U07-SRC-HELP04
% teaching: NEW-U07-VOC-tha-adverb-1
% teaching: NEW-U07-VOC-fisccynn-fish-kinds-source
% mdparagraph:P333:begin
\textbf{Then the great whales; all living and moving fish-kinds.} The first ðā after gescēop means “then”; the second is the article with miclan hwalas. Fisccynn is neuter accusative plural; libbendu and styriendlicu describe it. The source’s styriendlicu is supplied as an actual -u variant, distinct from the selected earlier styriendlic cell. The and joins the descriptions; it does not require a second newly invented noun meaning.
% mdparagraph:P333:end

\reserve{5\baselineskip}
% mdparagraph:P334:begin
\textbf{5. ðe ðā wæteru tugon forð on hiera hīwum.}
% mdparagraph:P334:end

% support: U07-SRC-HELP05
% teaching: NEW-U07-VOC-the-relative-which
% teaching: NEW-U07-VOC-hiw-kind-source
% teaching: NEW-U07-VOC-on-according-form-source
% teaching: NEW-U07-GRAM-U07-SOURCE-RELATIVE
% teaching: NEW-U07-GRAM-U07-SOURCE-CLASSIFICATION
% mdparagraph:P335:begin
\textbf{Which the waters brought forth according to their forms/\allowbreak{}kinds.} Supply the whole relative clause. The relative ðe is its object; ðā wæteru is its subject. Hīwum is dative plural of neuter hīw, “form; kind”. The phrase on hiera hīwum describes the classification of the creatures brought forth. Its local meaning is supplied rather than derived from the ordinary location sense of on.
% mdparagraph:P335:end

\reserve{5\baselineskip}
% mdparagraph:P336:begin
\textbf{6. æfter hiera cynne.}
% mdparagraph:P336:end

% support: U07-SRC-HELP06
% teaching: NEW-U07-VOC-aefter-according-to-source
% mdparagraph:P337:begin
\textbf{According to their kind.} Æfter governs dative cynne here. Supply this classificatory phrase with the flying kinds. It does not establish every use of æfter for independent composition.
% mdparagraph:P337:end

\reserve{5\baselineskip}
% mdparagraph:P338:begin
\textbf{7. God geseah ðā ðæt hit gōd wæs.}
% mdparagraph:P338:end

% support: U07-SRC-HELP07
% teaching: NEW-U07-VOC-thaet-conjunction-1
% teaching: NEW-U07-GRAM-U07-SOURCE-CONTENT
% mdparagraph:P339:begin
\textbf{God then saw that it was good.} The conjunction ðæt introduces the supplied content clause; its verb wæs follows the predicate gōd. Hit refers to the nonhuman situation described. Supply this whole clause relation and word order; recognizing its familiar indicative forms does not require general content-clause production.
% mdparagraph:P339:end

\reserve{5\baselineskip}
% mdparagraph:P340:begin
\textbf{8. And blētsode hīe, ðus cweðende.}
% mdparagraph:P340:end

% support: U07-SRC-HELP08
% teaching: NEW-U07-VOC-bletsian-bless-source
% teaching: NEW-U07-SRC-PART-FORMS
% teaching: NEW-U07-GRAM-U07-SOURCE-ACCOMPANYING-PARTICIPLE
% teaching: NEW-U07-GRAM-U07-SOURCE-SUBJECT-ELLIPSIS
% mdparagraph:P341:begin
\textbf{And [God] blessed them, saying thus.} God continues as the subject without being repeated; hīe is the accusative plural object. Blētsode is supplied as the third-singular preterite of plain blētsian. It is not the prefixed geblētsian of the productive tables. Cweðende refers to God and accompanies the blessing, introducing his words. This attendant participle construction is supplied even though the participle’s form is learned.
% mdparagraph:P341:end

\reserve{5\baselineskip}
% mdparagraph:P342:begin
\textbf{9. Weaxað … bēoð … gefyllað; ðā fuglas bēon gemanigfielde.}
% mdparagraph:P342:end

% support: U07-SRC-HELP09
% teaching: NEW-U07-VOC-gefyllan-fill-source
% mdparagraph:P343:begin
\textbf{Increase, be multiplied, and fill; let the birds be multiplied.} The created beings are first addressed with plural imperatives. The birds are then named as the third-person plural subject of a subjunctive wish. Gemanigfielde is the source’s nominative plural past participle in both clauses. Gefyllað means “fill” here; this lexical sense is supplied in addition to the earlier “finish, accomplish”.
% mdparagraph:P343:end

% mdparagraph:P344:begin
\textbf{Ðǣre sǣ wæteru} means “the waters of the sea”. Wæteru is the head noun and the thing to be filled. Sǣ is feminine genitive singular in this source phrase, and ðǣre agrees with it. This is not a prepositional location, nor is it the masculine genitive sǣs of the earlier selected paradigm. Keep the original gender and form in this supplied phrase.
% mdparagraph:P344:end

\reserve{5\baselineskip}
% mdparagraph:P345:begin
\textbf{10. And ðā wæs geworden ǣfen and męrgen se fīfta dæg.}
% mdparagraph:P345:end

% support: U07-SRC-HELP10
% teaching: NEW-U07-VOC-geweorthan-1
% teaching: NEW-U07-GRAM-U07-SOURCE-RESULTATIVE
% mdparagraph:P346:begin
\textbf{And then evening and morning had come to be, the fifth day.} Supply this resultative closing formula with intransitive geweorðan. It is not the passive of a transitive action: no accusative object has become the subject. Preserve the source spelling męrgen. The formula does not require an independently constructed rule for its agreement or a unique general clause analysis.
% mdparagraph:P346:end

% mdparagraph:P347:begin
\textbf{Source presentation.} The original words, spelling, quantity marks, endings, and punctuation are retained. Printed line wrapping is joined; the line-end divisions swim-/\allowbreak{}mendu and styriend-/\allowbreak{}licu are joined as words. The original note cues are moved to the list below. Curved Unicode quotation marks retain the source’s quotation boundaries; męrgen is not normalized to morgen.
% mdparagraph:P347:end

\reserve{7\baselineskip}
% mdparagraph:P348:begin
\subsection*{Cook’s original notes}
% mdparagraph:P348:end

% mdparagraph:P349:begin
The note numbers and section references below belong to Cook’s edition. The local help supplies their necessary interpretation without requiring the learner to consult those sections.
% mdparagraph:P349:end

% source-note: FB126-n4
% mdparagraph:P350:begin
\textbf{P.126, note 4 (ēac swilce).} Ēac swilce = etiam.
% mdparagraph:P350:end

% source-note: FB126-n5
% mdparagraph:P351:begin
\textbf{P.126, note 5 (forð).} Producant = tēon . . . forð.
% mdparagraph:P351:end

% source-note: FB126-n6
% mdparagraph:P352:begin
\textbf{P.126, note 6 (cucu).} See cucu.
% mdparagraph:P352:end

% source-note: FB126-n7
% mdparagraph:P353:begin
\textbf{P.126, note 7 (līfe).} See līf.
% mdparagraph:P353:end

% source-note: FB126-n8
% mdparagraph:P354:begin
\textbf{P.126, note 8 (flēogendu).} See flēogan, and 61.
% mdparagraph:P354:end

% source-note: FB126-n9
% mdparagraph:P355:begin
\textbf{P.126, note 9 (ðā (temporal, after gescēop)).} Adverb; see 84. 1.
% mdparagraph:P355:end

% source-note: FB126-n10
% mdparagraph:P356:begin
\textbf{P.126, note 10 (hwalas).} See hwæl.
% mdparagraph:P356:end

% source-note: FB126-n11
% mdparagraph:P357:begin
\textbf{P.126, note 11 (styriendlicu).} Lat. motabilem.
% mdparagraph:P357:end

% source-note: FB126-n12
% mdparagraph:P358:begin
\textbf{P.126, note 12 (ðe).} Acc.
% mdparagraph:P358:end

% source-note: FB126-n13
% mdparagraph:P359:begin
\textbf{P.126, note 13 (ðā (subject of tugon)).} Nom. plur.
% mdparagraph:P359:end

% source-note: FB126-n14
% mdparagraph:P360:begin
\textbf{P.126, note 14 (tugon).} See tēon.
% mdparagraph:P360:end

% source-note: FB126-n15
% mdparagraph:P361:begin
\textbf{P.126, note 15 (forð (after tugon)).} Tugon forð = producerunt.
% mdparagraph:P361:end

% source-note: FB126-n16
% mdparagraph:P362:begin
\textbf{P.126, note 16 (blētsode).} See blētsian, and 33.
% mdparagraph:P362:end

% source-note: FB126-n17
% mdparagraph:P363:begin
\textbf{P.126, note 17 (cweðende).} See cweðan.
% mdparagraph:P363:end

% source-note: FB126-n18
% mdparagraph:P364:begin
\textbf{P.126, note 18 (Weaxað).} See weaxan, and 24.
% mdparagraph:P364:end

% source-note: FB126-n19
% mdparagraph:P365:begin
\textbf{P.126, note 19 (gemanigfielde (first occurrence)).} Past part. in nom. plur.
% mdparagraph:P365:end

% source-note: FB126-n20
% mdparagraph:P366:begin
\textbf{P.126, note 20 (gefyllað).} See gefyllan.
% mdparagraph:P366:end

% source-note: FB126-n21
% mdparagraph:P367:begin
\textbf{P.126, note 21 (bēon).} See 193. a.
% mdparagraph:P367:end

\reserve{7\baselineskip}
% mdparagraph:P368:begin
\subsection*{Practice R. The original paragraph}
% mdparagraph:P368:end

% mdparagraph:P369:begin
Keep the printed local help available. Recognize the exact three participles and analyse taught forms within their supported clauses. Do not turn a case-syncretic spelling into a unique diagnosis.
% mdparagraph:P369:end

\reserve{5\baselineskip}
\par\begin{minipage}{\linewidth}\raggedright
% mdparagraph:P370:begin
\textbf{R1.}
% mdparagraph:P370:end

% item: NEW-U07-R01
% mdparagraph:P371:begin
Translate the complete fifth-day Creation paragraph with its printed local support.
% mdparagraph:P371:end

\end{minipage}\par\addvspace{4pt}

\par\begin{minipage}{\linewidth}\raggedright
% mdparagraph:P372:begin
\textbf{R2.}
% mdparagraph:P372:end

% item: NEW-U07-R02
% mdparagraph:P373:begin
Parse Tēon, Weaxað, bēoð, and bēon in their actual source clauses. Identify the subjects/\allowbreak{}addressees and explain the commands or wishes.
% mdparagraph:P373:end

\end{minipage}\par\addvspace{4pt}

\par\begin{minipage}{\linewidth}\raggedright
% mdparagraph:P374:begin
\textbf{R3.}
% mdparagraph:P374:end

% item: NEW-U07-R03
% mdparagraph:P375:begin
Find swimmendu, flēogendu, and libbendu in the paragraph. Give the noun each modifies and the case, number, and gender required at each occurrence.
% mdparagraph:P375:end

\end{minipage}\par\addvspace{4pt}

\par\begin{minipage}{\linewidth}\raggedright
% mdparagraph:P376:begin
\textbf{R4.}
% mdparagraph:P376:end

% item: NEW-U07-R04
% mdparagraph:P377:begin
Explain gemanigfielde in the blessing. Identify its two subjects and explain why the source’s plural ending is compatible with the verbal passive.
% mdparagraph:P377:end

\end{minipage}\par\addvspace{4pt}

\par\begin{minipage}{\linewidth}\raggedright
% mdparagraph:P378:begin
\textbf{R5.}
% mdparagraph:P378:end

% item: NEW-U07-R05
% mdparagraph:P379:begin
In ðǣre sǣ wæteru, identify the head noun and the dependent genitive. Translate the phrase and distinguish it from on ðǣm sǣ. Consult the local note on the source’s feminine sea-form.
% mdparagraph:P379:end

\end{minipage}\par\addvspace{4pt}

\par\begin{minipage}{\linewidth}\raggedright
% mdparagraph:P380:begin
\textbf{R6.}
% mdparagraph:P380:end

% item: NEW-U07-R06
% mdparagraph:P381:begin
Identify the limits of the relative clause with ðe using the supplied support. State what the waters brought forth, what God blessed, and which day ends the paragraph.
% mdparagraph:P381:end

\end{minipage}\par\addvspace{4pt}

\reserve{7\baselineskip}
% mdparagraph:P382:begin
\section*{11. Cumulative practice}
% mdparagraph:P382:end

\reserve{7\baselineskip}
% mdparagraph:P383:begin
\subsection*{Practice X. Combining the constructions}
% mdparagraph:P383:end

% mdparagraph:P384:begin
Preserve the earlier dative government, recipient constructions, purpose infinitive, and pronoun reference while using this unit’s mood, passive, and naming patterns.
% mdparagraph:P384:end

\reserve{5\baselineskip}
\par\begin{minipage}{\linewidth}\raggedright
% mdparagraph:P385:begin
\textbf{X1.}
% mdparagraph:P385:end

% item: NEW-U07-X01
% mdparagraph:P386:begin
A learner uses an accusative object after gīeman “rule over.” Explain the case required by this unit’s frame, then translate “The man rules over the birds.”
% mdparagraph:P386:end

\end{minipage}\par\addvspace{4pt}

\par\begin{minipage}{\linewidth}\raggedright
% mdparagraph:P387:begin
\textbf{X2.}
% mdparagraph:P387:end

% item: NEW-U07-X02
% mdparagraph:P388:begin
Translate “each tree” as a nominative subject and “all trees” as an accusative object. Explain the different number and agreement choices. Do not invent a plural meaning of ǣlc.
% mdparagraph:P388:end

\end{minipage}\par\addvspace{4pt}

\par\begin{minipage}{\linewidth}\raggedright
% mdparagraph:P389:begin
\textbf{X3.}
% mdparagraph:P389:end

% item: NEW-U07-X03
% mdparagraph:P390:begin
Translate “He can help the man, and he wants to give the birds food.” Identify the case governed by helpan and the two arguments of giefan.
% mdparagraph:P390:end

\end{minipage}\par\addvspace{4pt}

\par\begin{minipage}{\linewidth}\raggedright
% mdparagraph:P391:begin
\textbf{X4.}
% mdparagraph:P391:end

% item: NEW-U07-X04
% mdparagraph:P392:begin
Translate “I call the firmament Heaven,” then change only “I” to “he” and keep the naming relation. Identify both accusatives.
% mdparagraph:P392:end

\end{minipage}\par\addvspace{4pt}

\par\begin{minipage}{\linewidth}\raggedright
% mdparagraph:P393:begin
\textbf{X5.}
% mdparagraph:P393:end

% item: NEW-U07-X05
% mdparagraph:P394:begin
Give a second-plural command and a third-plural wish for “be multiplied.” Explain which word changes and what it tells the reader.
% mdparagraph:P394:end

\end{minipage}\par\addvspace{4pt}

\par\begin{minipage}{\linewidth}\raggedright
% mdparagraph:P395:begin
\textbf{X6.}
% mdparagraph:P395:end

% item: NEW-U07-X06
% mdparagraph:P396:begin
Translate: The companions came to help. The work was accomplished. Use the selected inflected purpose infinitive in the first sentence and a past passive with wesan in the second.
% mdparagraph:P396:end

\end{minipage}\par\addvspace{4pt}

\clearpage
% mdparagraph:P397:begin
\section*{12. Sources and selected forms}
% mdparagraph:P397:end

\begingroup\fontsize{10}{12.8}\selectfont
\setlength{\parskip}{1.5pt}
% mdparagraph:P398:begin
The principal source is Albert S. Cook, \textbf{A First Book in Old English: Grammar, Reader, Notes, and Vocabulary}, second edition, revised and enlarged (Ginn, 1895). Cook’s \textbf{Exercises in Old English} (Ginn, 1895) supplies I.4.18 and the complete I.6–7 prompts. They retain their original identifiers and wording despite being grouped in this grammatical unit.
% mdparagraph:P398:end

% mdparagraph:P399:begin
The full tables are edited teaching paradigms derived from the checked systems and selected stems. A completed cell is not a claim that the entire table occurs as a quotation in running prose. Joseph Wright and Elizabeth Mary Wright, \textbf{Old English Grammar}, second edition (Oxford, 1914), and R. D. Fulk, \textbf{An Introductory Grammar of Old English with an Anthology of Readings}, author-corrected open-access reprint (2020), corroborate the stated forms, variants, and bounded uses.
% mdparagraph:P399:end

% mdparagraph:P400:begin
All PDF pages are one-based. The selected Cook \emph{First Book} scan is printed page +16; Wright is +20; Fulk is +12. The \emph{Exercises} scan is printed page +8. The original Creation excerpt is Cook Reader I, p.126 /\allowbreak{} PDF p.142, editorial prose lines 5–14, with original notes 4–21. It is not an adaptation of the new account.
% mdparagraph:P400:end

% mdparagraph:P401:begin
\textbf{Subjunctive formation.} Cook §§93–95, pp.54–55 /\allowbreak{} PDF pp.70–71, supplies the basic stem and ending relations. The selected strong systems are in §§102–110, pp.59–67 /\allowbreak{} PDF pp.75–83; weak systems in §§113–124, pp.68–76 /\allowbreak{} PDF pp.84–92; preterite-presents and anomalous systems in §§126–141, pp.76–80 /\allowbreak{} PDF pp.92–96. Fulk §§45 and 88, pp.29 and 59 /\allowbreak{} PDF pp.41 and 71, corroborates the ending system; §123, pp.89–90 /\allowbreak{} PDF pp.101–102, distinguishes defective citation labels.
% mdparagraph:P401:end

% mdparagraph:P402:begin
\textbf{Mood and clauses.} Cook §§192–198, pp.103–105 /\allowbreak{} PDF pp.119–121, and Fulk §§44–45, pp.28–29 /\allowbreak{} PDF pp.40–41, support the elementary uses and the warning that gif does not automatically select subjunctive mood. The nonpast preterite example is a new composition with its unreal setting supplied; wider conditional syntax is deferred.
% mdparagraph:P402:end

% mdparagraph:P403:begin
\textbf{Participles and passive.} Cook §§55–62, pp.38–41 /\allowbreak{} PDF pp.54–57, gives the adjective and participial patterns; §179, pp.99–100 /\allowbreak{} PDF pp.115–116, and §§187–188, pp.101–102 /\allowbreak{} PDF pp.117–118, support the predicate and passive distinctions. Wright §§429–442, pp.211–218 /\allowbreak{} PDF pp.231–238, and Fulk §§54–56, pp.36–37 /\allowbreak{} PDF pp.48–49, corroborate agreement and syncope. The fuller stem spellings in the tables make the endings recoverable; the listed syncopated and plural variants remain valid.
% mdparagraph:P403:end

% mdparagraph:P404:begin
\textbf{Lexical frames.} Cook §156f, p.91 /\allowbreak{} PDF p.107, gives the genitive frame of gīeman; §173, p.97 /\allowbreak{} PDF p.113, gives the two-accusative naming relation. The new noun follows §47, p.32 /\allowbreak{} PDF p.48; ǣlc is identified in §89a, p.52 /\allowbreak{} PDF p.68, with its vocabulary entry on p.274 /\allowbreak{} PDF p.290. Spatial ofer with accusative and dative is explicit in Fulk §36, p.23 /\allowbreak{} PDF p.35; the glossary’s case mark is explained on p.251 /\allowbreak{} PDF p.263 and applied to ofer on p.297 /\allowbreak{} PDF p.309.
% mdparagraph:P404:end

% mdparagraph:P405:begin
The checked Cook vocabulary entries include gīeman (p.296 /\allowbreak{} PDF p.312), rīcsian (p.309 /\allowbreak{} PDF p.325), ge-manig-fieldan (p.293 /\allowbreak{} PDF p.309), tō-twǣman (p.316 /\allowbreak{} PDF p.332), wyrcean (p.324 /\allowbreak{} PDF p.340), and ge-weald (p.295 /\allowbreak{} PDF p.311). Dictionary compound hyphens are omitted in the learner headwords gemanigfieldan, tōtwǣman, and geweald; the source’s running spelling fisccynn is retained. No vowel length is added to gemanigfieldan.
% mdparagraph:P405:end

% mdparagraph:P406:begin
The additional lexical help follows Cook’s entries æfter (p.274 /\allowbreak{} PDF p.290), blētsian (p.279 /\allowbreak{} PDF p.295), cucu and cynn (pp.281–282 /\allowbreak{} PDF pp.297–298), ēac (p.283 /\allowbreak{} PDF p.299), fisc-cynn (p.287 /\allowbreak{} PDF p.303), forð (p.288 /\allowbreak{} PDF p.304), hīw (p.299 /\allowbreak{} PDF p.315), nū and ofer (p.307 /\allowbreak{} PDF p.323), swā (p.314 /\allowbreak{} PDF p.330), swilce and tēon (p.315 /\allowbreak{} PDF p.331), and ðus (p.319 /\allowbreak{} PDF p.335). These references support the particular meanings taught or supplied here, not every dictionary sense.
% mdparagraph:P406:end

% mdparagraph:P407:begin
The chapter’s eleven examples and “Binding the animals” are new compositions. The source’s relative and content clauses, accompanying participle, classification expressions, plain blētsode, and closing formula remain local help. Earlier instrumental and reference-only boundaries remain in force; complete morphology does not by itself teach every construction in which a form can occur.
% mdparagraph:P407:end

\endgroup
\end{document}
